\chapter{哈密顿\personFloer{}同调}
\label{ch:hamiltonian-floer}

\personFloer{}在 (Fl89b) 中为证明\personArnold{}猜想而引入了哈密顿\personFloer{}同调。为处理横截性
问题，他把讨论限制在单调辛流形的情形；这一概念也是他在同一篇论文中引入的。
\stfhIndex{I112}{Huofeier-Salameng}{霍费尔--萨拉蒙（Hofer--Salamon）}
随后，\personHofer{}与\personSalamon{} (HS95) 以及\personOno{} (On95) 把这个构造推广到了半正情形。

利用障碍丛的思想，在完全一般的情形下构造\personFloer{}复形，最早由\personKontsevich{}
(Kon95) 提出。\personFukaya{}与\personOno{} (FOn99)、\personLiu{}与\personGangTian{} (LT98) 以及\personYongbinRuan{} (Ru99)
实现了这一思想。按照\personFukaya{}与\personOno{}的术语，\termKuranishiStructure{}以及多值截面的使用，
对于获得上述横截性、进而证明基本模闭链的存在起着关键作用。

本节紧随 (Oh05c) 的叙述，把讨论限制在半正情形。本节末将简要说明如何利用\termKuranishiStructure{}把这一构造推广到一般情形。

尽管在半正情形可以取 $\mathbb Z$ 系数，我们仍将使用 $\mathbb Q$ 系数，因为在
使用\termKuranishiStructure{}的一般框架中需要后者。

\section{\personNovikov{} \personFloer{}链与\termNovikovRing{}}
\label{sec:novikov-floer-chains-ring}

设 $\phi=\phi_H^1\in\operatorname{Ham}(M,\omega)$ 非退化。对每个这样的
$H:S^1\times M\to\mathbb R$，我们知道 $\operatorname{Per}(H)$ 的基数有限。
考虑由 $\mathcal A_H$ 的临界点集
\[
 \operatorname{Crit}\mathcal A_H
 =\{[z,w]\in\widetilde{\mathcal L}_0(M)\mid z\in\operatorname{Per}(H)\}
\]
生成的自由 $\mathbb Q$-向量空间。为了正确定义\personFloer{}边界算子，需要关于元素
$[z,w]$ 的作用量 $\mathcal A_H([z,w])$ 所给出的实滤过，向下完备化这个向量空间。

\begin{definition}
\label{def:novikov-floer-chain}
考虑形式和
\begin{equation}
 \beta=\sum_{[z,w]\in\operatorname{Crit}\mathcal A_H}
 a_{[z,w]}[z,w],\qquad a_{[z,w]}\in\mathbb Q.
\label{eq:19-formal-floer-sum}
\end{equation}
\begin{enumerate}[label=(\arabic*)]
\stfhSubIndex{I142}{NuoweikefuFuluoerlian}{诺维科夫弗洛尔链（Novikov Floer chain）}{Zhiji}{支集（support）}
\item 定义 $\beta$ 的\emph{支集}为
\[
 \operatorname{supp}(\beta):=
 \{[z,w]\in\operatorname{Crit}\mathcal A_H\mid
 a_{[z,w]}\neq0\text{，在和式~\eqref{eq:19-formal-floer-sum} 中}\}.
\]
\stfhSubIndex{I140}{NuoweikefuFuluoerlian}{诺维科夫弗洛尔链（Novikov Floer chain）}{Shengchengyuan}{生成元（generator）}
\item 称满足 $a_{[z,w]}\neq0$ 的 $[z,w]$ 为和式 $\beta$ 的\emph{生成元}，并记作
\[
 [z,w]\in\operatorname{supp}\beta.
\]
在这种情形下，也称 $[z,w]$ 对 $\beta$ \emph{有贡献}。
\item 若形式和 $\beta$ 对任意 $\lambda\in\mathbb R$ 都满足
\begin{equation}
 \#\bigl(\operatorname{supp}(\beta)\cap
 \{[z,w]\mid\mathcal A_H([z,w])\geq\lambda\}\bigr)<\infty,
\label{eq:19-novikov-finiteness}
\end{equation}
\stfhIndex{I139}{NuoweikefuFuluoerlian}{诺维科夫弗洛尔链（Novikov Floer chain）}
则称 $\beta$ 为\emph{\personNovikov{} \personFloer{}链}（简称\emph{\personFloer{}链}）。称此条件为\termNovikovFinitenessCondition{}，并以 $CF_*(H)$ 表示\personFloer{}链的集合。
\end{enumerate}
\end{definition}

注意，按此定义，$CF_*(H)$ 是一个 $\mathbb Q$-向量空间，一般总是无限维的；只有当
$(M,\omega)$ 辛非球面时，它才是有限维的。给每个生成元
$[z,w]\in\operatorname{Crit}\mathcal A_H$ 配上\personConley{}--\personZehnder{}指标
$\mu_H([z,w])$，便赋予 $CF_*(H)$ 分次向量空间结构。

现在考虑\personFloer{}链
\[
 \beta=\sum a_{[z,w]}[z,w],\qquad a_{[z,w]}\in\mathbb Q.
\]
下面的概念是借助\personFloer{}同调建立极小极大临界点理论的关键。

\stfhSubIndex{I141}{NuoweikefuFuluoerlian}{诺维科夫弗洛尔链（Novikov Floer chain）}{Feng}{峰（peak）}
\begin{definition}
\label{def:floer-chain-level-peak}
设 $0\neq\beta\in CF_*(H)$ 是\personFloer{}链。定义链 $\beta$ 的\emph{层级}为
\[
 \lambda_H(\beta)=\max_{[z,w]}
 \{\mathcal A_H([z,w])\mid[z,w]\in\operatorname{supp}(\beta)\},
\]
并令 $\lambda_H(0)=-\infty$。若生成元 $[z,w]\in\beta$ 满足
$\mathcal A_H([z,w])=\lambda_H(\beta)$，则称它为 $\beta$ 的一个\emph{峰}，并记作
$\operatorname{peak}(\beta)$。
\end{definition}

我们强调，正是定义 19.1.1 中的\personNovikov{}有限性条件~\eqref{eq:19-novikov-finiteness}
保证了 $\lambda_H(\beta)$ 良定义。下述引理说明\personNovikov{}覆盖空间的定义是最优的。

\begin{lemma}
\label{lem:unique-peak-fixed-orbit}
设 $\beta\neq0$ 是齐次\personFloer{}链。则在一个固定周期轨道上，$\beta$ 的峰是唯一的。
\end{lemma}

\begin{proof}
设 $[z,w]$ 与 $[z,w']$ 是 $\beta$ 的两个这样的峰。于是
\[
 \mathcal A_H([z,w])=\lambda_H(\beta)=\mathcal A_H([z,w']),
\]
从而 $\omega([w])=\omega([w'])$。由齐次性假设还有
\[
 \mu_H([z,w])=\mu_H([z,w']).
\]
根据 $\Gamma$-等价类的定义，$[z,w]=[z,w']$，证明完毕。
\end{proof}

到目前为止，我们已把 $CF_*(H)$ 定义为以
$\operatorname{Crit}\mathcal A_H$ 为生成集的 $\mathbb Z$-分次
$\mathbb Q$-向量空间。现在说明如何像 (Fl89b)、(HS95) 那样，把 $CF(H)$ 描述成\emph{\personNovikov{}环}上的模。

回忆 $\Gamma$ 是交换群。考虑由有限和
\[
 R=\sum_{i=1}^{k}r_iq^{A_i}\in\mathbb Q[\Gamma]
\]
组成的群环 $\mathbb Q[\Gamma]$，并定义其支集为
\[
 \operatorname{supp}R=
 \{A\in\Gamma\mid A=A_i\text{，且在此和式中 }r_i\neq0\}.
\]
回忆赋值 $\nu:\Gamma\to\mathbb R$ 与次数映射 $d:\Gamma\to\mathbb Z$
\footnote{译者注：原书此处把次数映射的陪域印作 $\mathbb R$；但此前式 (18.2.13)
已定义 $d:\Gamma\to\mathbb Z$、$d(A)=2c_1(A)$，译文依该定义改正。}
。
赋值 $\nu:\mathbb Q[\Gamma]\to\mathbb R$ 定义为
\[
 \nu(R)=\nu\left(\sum_{i=1}^{k}r_iq^{A_i}\right)
 :=\min\{\omega(A)\mid A\in\operatorname{supp}R\}.
\]
它满足如下非阿基米德三角不等式：
\begin{equation}
 \nu(R_1+R_2)\geq\min\{\nu(R_1),\nu(R_2)\},
\label{eq:19-nonarch-triangle}
\end{equation}
因而在 $\mathbb Q[\Gamma]$ 上诱导由下列距离给出的自然度量拓扑
\footnote{译者注：原书在此写 $e^{\nu(R_1-R_2)}$。若赋值趋于 $+\infty$ 表示\personNovikov{}级数收敛，则通常相应距离应含负指数；原式与紧接着按此赋值作“向上完备化”
的约定表面上不相容。信息不足以唯一判定作者的符号规范，故正文照录。}
：
\[
 d(R_1,R_2):=e^{\nu(R_1-R_2)}.
\]

\begin{definition}
\label{def:novikov-ring-upward-completion}
\personNovikov{}环是 $\mathbb Q[\Gamma]$ 关于赋值
$\nu:\mathbb Q[\Gamma]\to\mathbb R$ 的\emph{向上}完备化
$\mathbb Q[\Gamma\rangle$，记作 $\Lambda_\omega=\Lambda_\omega^\uparrow$。
\end{definition}

更具体地，
\[
 \Lambda_\omega=
 \left\{\sum_{A\in\Gamma}r_Aq^A\ \middle|\
 \forall\lambda\in\mathbb R,
 \#\{A\in\Gamma\mid r_A\neq0,\ \omega(A)<\lambda\}<\infty\right\}.
\]
在 $\operatorname{Crit}\mathcal A_H$ 上定义“粘合一个球面”的 $\Gamma$-作用
\[
 [z,w]\longmapsto[z,w\#(-A)],
\]
进而以和 $\Lambda_\omega$ 的卷积积在 $CF(H)$ 上诱导 $\Lambda_\omega$-模结构：
\begin{equation}
\begin{aligned}
 &\left(\sum_{A\in\Gamma}r_Aq^A\right)
 \cdot\left(\sum a_{[z,w]}[z,w]\right)\\
 &\quad:=\sum_{[z,w']}
 \left(\sum_{\substack{A,[z,w]\,;\\{}[z,(-A)\#w]=[z,w']}}
 r_Aa_{[z,(-A)\#w]}\right)[z,w'].
\end{aligned}
\label{eq:19-convolution-action}
\end{equation}
按照这个定义，$q^A[z,w]=[z,(-A)\#w]$。换言之，因为这里讨论的是同调而非上同调，
我们改变了符号。

以下我们力求始终以 $CF(H)$ 表示 $\Lambda_\omega$-模，以 $CF_*(H)$ 表示分次
$\mathbb Q$-向量空间。作用量泛函在 $CF_*(H)$ 上给出自然滤过：给定任意
$\lambda\in\mathbb R\setminus\operatorname{Spec}(H)$，定义
\[
 CF_*^\lambda(H)=\{\alpha\in CF_*(H)\mid\lambda_H(\alpha)\leq\lambda\},
\]
并以
\[
 i_\lambda:CF_*^\lambda(H)\longrightarrow CF_*(H)
\]
表示自然包含同态。

稍后会看到，把 $q^A=T^{\omega(A)}e^{c_1(A)}$ 中的形式变量 $T$ 与 $e$ 分开有时
很有用。记 $\Lambda_\omega^{(0)}$ 为 $\Lambda_\omega$ 的零次部分，即满足\personNovikov{}有限性条件的形式级数
\[
 \sum_{\substack{A\in\Gamma\\c_1(A)=0}}a_Aq^A
 =\sum_{\substack{A\in\Gamma\\c_1(A)=0}}a_AT^{\omega(A)}
\]
的集合。显然，我们需要验证下述结论。

\begin{proposition}
\label{prop:degree-zero-novikov-field}
$\Lambda_\omega^{(0)}$ 是域，并且
\[
 \Lambda_\omega=\Lambda_\omega^{(0)}[e,e^{-1}],
\]
即 $\Lambda_\omega$ 是域 $\Lambda_\omega^{(0)}$ 上关于 $e,e^{-1}$ 的多项式环。
\end{proposition}

\begin{exercise}
\label{ex:prove-degree-zero-novikov-field}
证明此命题。
\end{exercise}

\section{\personFloer{}边界映射的定义}
\label{sec:def-hamiltonian-floer-boundary}

设 $H$ 是非退化的一周期哈密顿量，$J$ 是相容殆复结构的一周期族。我们先定义\personFloer{}边界映射，并指出定义这一对的\personFloer{}同调 $HF_*(H,J)$ 所需的横截性条件。
在 $(M,\omega)$ 的半正性假设下，我们将完整讨论这个横截性问题。

\begin{definition}
\label{def:relative-cylinder-homotopy-class}
设 $z,z'\in\operatorname{Per}(H)$。以 $\pi_2(z,z')$ 表示相对于边界
\[
 u(0,t)=z(t),\qquad u(1,t)=z'(t)
\]
的光滑映射
\[
 u:[0,1]\times S^1=:T\longrightarrow M
\]
的同伦类集合，并以 $[u]\in\pi_2(z,z')$ 表示 $u$ 的同伦类。
\end{definition}

类似地，以 $\pi_2(z)$ 表示映射
\[
 w:D^2\longrightarrow M;\qquad w|_{\partial D^2}=z
\]
的相对同伦类集合。注意，$\pi_2(M)$ 在 $\pi_2(z)$ 和 $\pi_2(z,z')$ 上有自然的
“粘合一个球面”作用（模去 $\pi_1(M)$ 的作用）。此外，对
$C\in\pi_2(z,z')$，存在由粘合映射 $w\mapsto w\#u$ 诱导的自然映射
\[
 (\,\cdot\,)\#C:\pi_2(z)\longrightarrow\pi_2(z').
\]
具体地说，我们固定在 $D^2$ 的极坐标 $(r,\theta)$ 中按下式定义
$w\#u:D^2\to M$：
\begin{equation}
 (w\#u)(r,\theta)=
 \begin{cases}
  w(2r,\theta),&0\leq r\leq\frac12,\\
  u(2r-1,\theta),&\frac12\leq r\leq1.
 \end{cases}
\label{eq:19-disc-cylinder-gluing}
\end{equation}
还有自然的粘合映射
\[
 \pi_2(z_0,z_1)\times\pi_2(z_1,z_2)\longrightarrow\pi_2(z_0,z_2),
 \qquad (u_1,u_2)\longmapsto u_1\#u_2.
\]

\stfhIndex{I167}{XiangduiKangli-Zengdeerzhibiao}{相对康利--曾德尔指标（relative Conley--Zehnder index）}
现在定义\personConley{}--\personZehnder{}指标
$\mu_H:\pi_2(z,z')\to\mathbb Z$。只要 $z_0,z_1$ 位于 $\Omega(M)$ 的同一连通分支，
这个定义并不要求它们可缩。

令 $T=[0,1]\times S^1$。对给定映射 $u:T\to M$，选取一个辛平凡化
\[
 \Phi:u^*TM\longrightarrow T\times\mathbb R^{2n}.
\]
要求它满足标记条件
\begin{equation}
 \Phi(t,1)\equiv\Phi(0,1),\qquad t\in[0,1].
\label{eq:19-trivialization-marking}
\end{equation}
我们知道，对某些 $p_0,p_1\in\operatorname{Fix}\phi_H^1$，有
$z_0(t)=\phi_H^t(p_0)$、$z_1(t)=\phi_H^t(p_1)$。于是有两个映射
\[
 \alpha_{\Phi,i}:[0,1]\longrightarrow\operatorname{Sp}(2n),\qquad i=0,1,
\]
使得对 $v\in\mathbb R^{2n}$ 和 $t\in[0,1]$，
\[
 \Phi\circ d\phi_H^t(p_i)\circ\Phi^{-1}(i,t,v)
 =(i,t,\alpha_{\Phi,i}(t)v).
\]
由条件~\eqref{eq:19-trivialization-marking} 及 $H$ 的非退化性，对 $i=0,1$ 都有
\[
 \alpha_{\Phi,i}(0)=\operatorname{id},\qquad
 \alpha_{\Phi,i}(1)\in\operatorname{Sp}^*(2n),
\]
因而\personConley{}--\personZehnder{}指标 $\mu_{CZ}(\alpha_{\Phi,i})$ 有定义。

\begin{definition}
\label{def:relative-conley-zehnder-index}
对每个 $C\in\pi_2(z,z')$，定义其相对\personConley{}--\personZehnder{}指标为
\begin{equation}
 \mu_H(z,z';C):=\mu_{CZ}(\alpha_{\Phi,0})-
 \mu_{CZ}(\alpha_{\Phi,1}).
\label{eq:19-relative-cz-index}
\end{equation}
边界条件不会引起混淆时，也写作 $\mu_H(C)$。
\end{definition}

\begin{exercise}
\label{ex:relative-cz-trivialization-independent}
证明此定义与满足~\eqref{eq:19-trivialization-marking} 的平凡化 $\Phi$ 无关。
\end{exercise}

下面的命题把这个相对指标与
$[z,w]\in\operatorname{Crit}\mathcal A_H$ 的绝对\personConley{}--\personZehnder{}指标
$\mu_H([z,w])$ 联系起来。

\begin{proposition}
\label{prop:relative-absolute-cz-relation}
对类 $C\in\pi_2(z,z')$ 的\emph{一个（因而任意一个）代表}
$u:[0,1]\times S^1\to M$
\footnote{译者注：原书在代表映射的类型中印作
$u:[0,1]\times S^1\times M$，缺少映射箭头；结合本节定义及随后证明，译文补为
$u:[0,1]\times S^1\to M$。}
，有
\begin{equation}
 \mu_H(z,z';C)=\mu_H([z,w])-\mu_H([z',w\#u]).
\label{eq:19-relative-absolute-cz}
\end{equation}
\end{proposition}

\begin{proof}
选取平凡化
\[
 \Psi_w:w^*TM\longrightarrow D^2\times\mathbb R^{2n}
\]
以及满足~\eqref{eq:19-trivialization-marking} 和
\begin{equation}
 \Psi_w|_{\partial D^2}=\Phi|_{\{0\}\times S^1}
\label{eq:19-compatible-glued-trivializations}
\end{equation}
的平凡化 $\Phi:u^*TM\to T\times\mathbb R^{2n}$；这里把
$\partial D^2\cong\{0\}\times S^1$。粘合 $\Psi_w\#\Phi$ 自然诱导
$(w\#u)^*TM$ 的平凡化。由构造和条件~\eqref{eq:19-compatible-glued-trivializations}，
\[
 \mu_H(\alpha_{\Psi_w\#\Phi})=\mu_H(\alpha_{\Phi,1}),
 \qquad
 \mu_H(\alpha_{\Psi_w})=\mu_H(\alpha_{\Phi,0}).
\]
因此
\[
\begin{aligned}
 \mu_H([z',w\#u])
 &=\mu_H(\alpha_{\Psi_w\#\Phi})\\
 &=\mu_H(\alpha_{\Phi,1})-\mu_H(\alpha_{\Phi,0})
   +\mu_H(\alpha_{\Phi,0})\\
 &=-\mu_H(z,z';[u])+\mu_H([z,w]),
\end{aligned}
\]
证明完毕。
\footnote{译者注：原书本计算首项漏印右方括号，写成
$\mu_H([z',w\#u]$；译文补齐。}
\end{proof}

现在以
\[
 \widetilde{\mathcal M}(H,J;z,z';C)
\]
表示满足渐近条件与同伦条件
\begin{equation}
 u(-\infty)=z,\qquad u(\infty)=z';\qquad [u]=C
\label{eq:19-floer-asymptotic-homotopy}
\end{equation}
的方程 (18.5.41) 的有限能量解集，并令
\[
 \mathcal M(H,J;z,z';C)
 =\widetilde{\mathcal M}(H,J;z,z';C)/\mathbb R.
\]
这里要指出，虽然 $u$ 定义在 $\mathbb R\times S^1$ 上，但由于方程 (18.5.41) 的
解具有指数衰减性质（回忆我们假设 $H$ 非退化），它可紧化为满足相应边界条件
\[
 \bar u(0)=z,\qquad\bar u(1)=z'
\]
的连续映射 $\bar u:[0,1]\times S^1\to M$。称 $\bar u$ 为 $u$ 的\emph{紧化映射}。
略用记号，也以 $[u]$ 表示紧化映射的类 $[\bar u]\in\pi_2(z,z')$。

下面是\personFredholm{}指标的基本公式。（证明见 (Schw95)。）

\begin{proposition}
\label{prop:closed-floer-index-formula}
设 $u\in\mathcal M(H,J;z,z';C)$，则
\begin{equation}
 \operatorname{Index}u=\mu_H(z,z';C).
\label{eq:19-closed-floer-index}
\end{equation}
特别地，若 $J$ 是 $H$-正则的，则
\[
 \dim\widetilde{\mathcal M}(H,J;z,z';C)=\mu_H(z,z';C).
\]
\end{proposition}

\begin{exercise}
\label{ex:closed-string-index-formula}
把命题 15.3.1 的证明转写到闭弦情形，并证明此命题。
\end{exercise}

\personFloer{}边界映射
\[
 \partial_{(H,J)}:CF_{k+1}(H)\longrightarrow CF_k(H)
\]
通过计数满足 $\mu_H(z,z';C)=1$ 的三元组 $(z,z';C)$ 所对应的
$\mathcal M(H,J;z,z';C)$ 中的元素来定义，条件如下。

\begin{definition}[边界映射]
\label{def:hamiltonian-floer-boundary-regularity}
设 $H$ 非退化。假设 $J$ 满足：
\begin{enumerate}[label=(\arabic*)]
\item 对任意满足
\[
 \mu_H(z_0,z_1;C)
 =\mu_H([z_0,w_0])-\mu_H([z_1,w_0\#C])\leq0
\]
的对子 $z_0,z_1\in\operatorname{Per}(H)$，除非 $z_0=z_1$ 且 $C=0$，否则
$\widetilde{\mathcal M}(H,J;z_0,z_1;C)=\varnothing$。当 $z_0=z_1$ 且 $C=0$ 时，
仅有定常解，即对所有 $\tau\in\mathbb R$ 都有 $u(\tau)\equiv z_0=z_1$。
\item 对任意 $z_0,z_1\in\operatorname{Per}(H)$ 以及满足
\[
 \mu_H(z_0,z_1;C)=1
\]
的同伦类 $C\in\pi_2(z_0,z_1)$，空间
$\widetilde{\mathcal M}(H,J;z_0,z_1;C)$ 横截，且商
\[
 \mathcal M(H,J;z_0,z_1;C)
 =\widetilde{\mathcal M}(H,J;z_0,z_1;C)/\mathbb R
\]
紧致，因而是有限集。记
\[
 n(H,J;z_0,z_1;C)=\#\mathcal M(H,J;z_0,z_1;C)
\]
为此空间中元素的代数计数。若三元组 $(z,z';C)$ 不满足指标要求，则约定
\[
 n(H,J;z_0,z_1;C)=0.
\]
\item 对任意 $z_0,z_2\in\operatorname{Per}(H)$ 和满足
\[
 \mu_H(z_0,z_2;C)=2
\]
的 $C\in\pi_2(z_0,z_2)$，$\mathcal M(H,J;z_0,z_2;C)$ 可紧化为光滑一维流形，
其边界由所有可能的断裂流线
\[
 [u_1]\#_\infty[u_2]
\]
组成，其中对某个 $y\in\operatorname{Per}(H)$ 以及
$C_1\in\pi_2(z_0,y)$、$C_2\in\pi_2(y,z_2)$，
\[
 u_1\in\widetilde{\mathcal M}(H,J;z_0,y;C_1),\qquad
 u_2\in\widetilde{\mathcal M}(H,J;y,z_2;C_2),
\]
并且
\[
 C_1\#C_2=C,\qquad [u_1]=C_1,\quad[u_2]=C_2,
\]
以及
\[
 \mu_H(z_0,y;C_1)=\mu_H(y,z_2;C_2)=1.
\]
这里以 $[u]$ 表示由 $u$ 代表的等价类。
\end{enumerate}
\stfhIndex{I089}{Fuluoerzhengze}{弗洛尔正则（Floer regular）}
称任何这样的 $J$ 为 \emph{$H$-正则}，并称这样的对子 $(H,J)$ 为\emph{\personFloer{}正则}。
\end{definition}

这个定义的要点是：对\personFloer{}正则对子 $(H,J)$，\personFloer{}边界映射
\[
 \partial=\partial_{(H,J)}:CF_*(H)\longrightarrow CF_*(H)
\]
有定义且满足 $\partial\partial=0$，于是可以取其同调。由于 $CF_*(H)$ 由
$[z,w]\in\operatorname{Crit}\mathcal A_H$ 生成，我们需要更详细地说明这个构造。

给定 $[z,w]$。对每个 $[z',w']\in\operatorname{Crit}\mathcal A_H$，考虑所有满足
\begin{equation}
 [z',w\#C]=[z',w'],\qquad\mu_H(z,z';C)=1
\label{eq:19-boundary-class-condition}
\end{equation}
的 $C\in\pi_2(z,z')$，并定义模空间
\[
 \mathcal M(H,J;[z,w],[z',w'])
 :=\bigcup_C\{\mathcal M(H,J;z,z';C)\mid
 C\text{ 满足~\eqref{eq:19-boundary-class-condition}}\}.
\]
按 $\Gamma$-覆盖空间 $\widetilde{\mathcal L}_0(M)$ 的定义，可能有不止一个
$C\in\pi_2(z^-,z^+)$ 满足~\eqref{eq:19-boundary-class-condition}。不过，\personGromov{}--\personFloer{}紧性给出如下有限性。

\begin{lemma}
\label{lem:finitely-many-boundary-classes}
上述并是有限并。换言之，对任意给定对子 $([z,w],[z',w'])$，只有有限多个
$C\in\pi_2(z,z')$ 同时满足~\eqref{eq:19-boundary-class-condition} 和
$\mathcal M(H,J;z,z';C)\neq\varnothing$。
\end{lemma}

\begin{proof}
反设存在无穷序列
$u_i\in\mathcal M(H,J;[z,w],[z',w'])$，且 $[u_i]\in\pi_2(z,z')$ 两两不同。
首先注意
\[
 E_{(H,J)}(u_i)=\mathcal A_H([z,w])-\mathcal A_H([z',w'])
\]
与 $i$ 无关。由\personGromov{}--\personFloer{}紧性，可取仍记为 $u_i$ 的子列，使其收敛到
\[
 u_\infty=u^1\#u^2\#\cdots\#u^N,
\]
其中每个 $u^k$ 形如
\[
 u^k=u_{k,0}+\sum_{\ell=1}^{L_k}v_\ell^k.
\]
这里，$\{u_{k,0}\}_k$ 构成一条从 $z$ 到 $z'$ 的连通\personFloer{}流线链，
$v_\ell^k$ 是附着于 $u_{k,0}$ 的球面气泡；每个 $u_{k,0}$ 是连接周期轨道
$z_{k-1},z_k$ 的光滑\personFloer{}流线，其中 $z_0=z$、$z_N=z'$。

此外，\personGromov{}--\personFloer{}紧性还蕴含：对所有充分大的 $i$，有与 $i$ 无关的等式
\footnote{译者注：原书此处只写
$[u_i]=[u_\infty]+\sum_{\ell=1}^{L_k}[v_\ell^k]$，其中 $k$ 未被量化；而上文
$u_\infty$ 含全部 $N$ 个分支。译文据该分解补出对 $k$ 的求和。}
\[
 [u_i]=[u_\infty]+\sum_{k=1}^{N}\sum_{\ell=1}^{L_k}[v_\ell^k].
\]
这与 $[u_i]\in\pi_2(z,z')$ 两两不同的假设矛盾。
\end{proof}

现在把 $u$ 看成覆盖空间 $\widetilde{\mathcal L}_0(M)$ 中的道路，把
$u\in\mathcal M(H,J;[z,w],[z',w'])$ 的渐近条件写成
\begin{equation}
 u(-\infty)=[z,w],\qquad u(\infty)=[z',w'].
\label{eq:19-lifted-asymptotic-condition}
\end{equation}
\personFloer{}边界映射
$\partial=\partial_{(H,J)}:CF_*(H)\to CF_*(H)$ 由矩阵系数
\[
 \langle\partial([z,w]),[z',w']\rangle
 :=\sum_C n_{(H,J)}(z,z';C)
\]
定义，其中 $C$ 如~\eqref{eq:19-boundary-class-condition}，且 $[z,w]$ 与 $[z',w']$
的\personConley{}--\personZehnder{}指标满足
\[
 \mu_H([z,w])-\mu_H([z',w'])=\mu(z,z';C)=1.
\]
由定义 19.2.7 的定义与假设，$\partial=\partial_{(H,J)}$ 的次数为 $-1$，且
$\partial\circ\partial=0$。

\stfhIndex{I084}{Fuluoerbilian}{弗洛尔闭链（Floer cycle）}
\begin{definition}
\label{def:floer-cycle-boundary-homology}
若\personFloer{}链 $\beta\in CF(H)$ 满足 $\partial\beta=0$，即
$\beta\in\ker\partial_{(H,J)}$，则称 $\beta$ 是 $(H,J)$ 的\emph{\personFloer{}闭链}；若
$\beta\in\operatorname{Image}\partial_{(H,J)}$，则称它为\emph{\personFloer{}边界}。若
$\beta'-\beta$ 是边界，则称两个\personFloer{}链 $\beta,\beta'$ \emph{同调}。
\end{definition}

定义 $(H,J)$ 的\personFloer{}同调为
\[
 HF_*(H,J):=\ker\partial_{(H,J)}/\operatorname{Image}\partial_{(H,J)}.
\]
它既可视为分次 $\mathbb Q$-向量空间，也可视为 $\Lambda_\omega$-模。

\stfhIndex{I082}{Fuluoerlianyingshe}{弗洛尔链映射（Floer chain map）}
\section{\personFloer{}链映射的定义}
\label{sec:def-floer-chain-map}

设给定一族 $(\mathcal H,\mathfrak j)$，其中
$\mathcal H=\{H^s\}_{0\leq s\leq1}$、
$\mathfrak j=\{J^s\}_{0\leq s\leq1}$，以及拉长函数
$\rho:\mathbb R\to[0,1]$。对每个这样的对子 $(\mathcal H,\mathfrak j)$ 与
$\rho$，$(\mathcal H,\mathfrak j;\rho)$ 携带\emph{非自治}方程 (18.5.42)，其渐近条件为
\begin{equation}
 u(-\infty)=z_0,\qquad u(\infty)=z_1,
\label{eq:19-chain-map-asymptotics}
\end{equation}
同伦条件为固定的 $[u]=C\in\pi_2(z_0,z_1)$。以
\[
 \mathcal M((\mathcal H,\mathfrak j;\rho);z_0,z_1;C)
\]
表示满足这些渐近边界条件与同伦条件的方程 (18.5.44) 的有限能量解集。

链同态
\[
 h_{\mathcal H}=h_{(\mathcal H,\mathfrak j;\rho)}:
 CF_*(H^0)\longrightarrow CF_*(H^1)
\]
通过遍历周期轨道对子 $z_0,z_1$ 与同伦类 $C$，考察这些模空间
$\mathcal M((\mathcal H,\mathfrak j;\rho);z_0,z_1;C)$ 来定义。容易调整定义 19.2.2
中的\personConley{}--\personZehnder{}指标定义，以得到链映射所需的指标。

\begin{exercise}
\label{ex:cz-index-continuation-triple}
调整定义 19.2.2，为连接 $H^0$ 与 $H^1$ 的同伦 $\mathcal H$ 所对应的三元组
$(z_0,z_1;C)$ 定义\personConley{}--\personZehnder{}指标，其中
$z_0\in\operatorname{Per}(H^0)$、$z_1\in\operatorname{Per}(H^1)$ 且
$C\in\pi_2(z_0,z_1)$。
\footnote{译者注：原书把 $z_1$ 与 $\operatorname{Per}(H^1)$ 的关系印作子集符号
$z_1\subset\operatorname{Per}(H^1)$；因 $z_1$ 是单条周期轨道，译文改为成员符号。}
\end{exercise}

\begin{definition}[链映射]
\label{def:floer-regular-continuation-data}
若下列条件成立，则称 $(\mathcal H^\rho,\mathfrak j^\rho)$ \emph{\personFloer{}正则}。
\begin{enumerate}[label=(\arabic*)]
\item 若 $\mu_{\mathcal H^\rho}(z_0,z_1;C)\leq-1$，则
\[
 \mathcal M(\mathcal H^\rho,\mathfrak j^\rho;z_0,z_1;C)=\varnothing.
\]
\item 对任意满足
\[
 \mu_{\mathcal H^\rho}(z_0,z_1;C)=0
\]
的 $z_0\in\operatorname{Per}(H^0)$ 和 $z_1\in\operatorname{Per}(H^1)$，
$\mathcal M(\mathcal H^\rho,\mathfrak j^\rho;z_0,z_1;C)$ 横截且紧致，因而为有限集。
记
\[
 n(\mathcal H^\rho,\mathfrak j^\rho;z_0,z_1;C)
 :=\#\mathcal M(\mathcal H^\rho,\mathfrak j^\rho;z_0,z_1;C)
\]
为其中元素的代数计数；其他情形令
$n(\mathcal H^\rho,\mathfrak j^\rho;z_0,z_1;C)=0$。
\item 对任意满足
\[
 \mu_{\mathcal H^\rho}(z_0,z_2;C)=1
\]
的 $z_0\in\operatorname{Per}(H^0)$ 和 $z_2\in\operatorname{Per}(H^1)$，
$\mathcal M(\mathcal H^\rho,\mathfrak j^\rho;z_0,z_2;C)$ 横截，并可紧化为光滑
一维流形，其边界由断裂流线 $u_1\#_\infty u_2$ 的全体组成；其中或者
\[
\begin{gathered}
 (u_1,u_2)\in
 \mathcal M(\mathcal H^\rho,\mathfrak j^\rho;z_0,y;C_1)
 \times\mathcal M(H^1,J^1;y,z_2;C_2),\\
 \mu_{\mathcal H^\rho}(z_0,y;C_1)=0,\qquad
 \mu_H(y,z_2;C_2)=1,
\end{gathered}
\]
或者
\[
\begin{gathered}
 (u_1,u_2)\in
 \mathcal M(H^0,J^0;z_0,y;C_1)
 \times\mathcal M(\mathcal H^\rho,\mathfrak j^\rho;y,z_2;C_2),\\
 \mu_H(z_0,y;C_1)=1,\qquad
 \mu_{\mathcal H^\rho}(y,z_2;C_2)=0,
\end{gathered}
\]
并且前一种断裂中 $y\in\operatorname{Per}(H^1)$，后一种断裂中
$y\in\operatorname{Per}(H^0)$；对所有相应的
$C_1\in\pi_2(z_0,y)$、$C_2\in\pi_2(y,z_2)$，有 $C_1\#C_2=C$
\footnote{译者注：原书本定义记号多处漂移：第 (1)--(2) 项把标题中的
$(\mathcal H^\rho,\mathfrak j^\rho)$ 写成 $(\mathcal H_{\mathbb R},\mathfrak j_{\mathbb R})$（下标为黑板体 $\mathbb R$）；第 (3) 项
又写未拉长的 $(H,J)$，并在第一类断裂的后一因子中把同伦类写成上标 $C^2$、
在第二类断裂的后一因子及同伦类中把 $C_2$ 误印为 $C_1$；第二类断裂的指标行
还把 $\mathcal H^\rho$ 与 $H$ 的槽位互错（印作
$\mu_{\mathcal H^\rho}(z_0,y;C_1)=1$、$\mu_H(y,z_2;C_2)=0$）。源书又通篇以
$\operatorname{Per}(H_0)$、$\operatorname{Per}(H_1)$ 与
$\operatorname{Crit}\mathcal A_{H_0}$、$\operatorname{Crit}\mathcal A_{H_1}$
的下标形式书写 $H^0,H^1$，并写 $y\in\operatorname{Per}(H)$，而同伦中没有单个
$H$。译文依本定义标题、紧邻的指标式、$C_1\#C_2=C$ 以及第一类断裂自身的
配对模式，统一为上标记号并配以相容的指标。}
。
\end{enumerate}
\end{definition}

对每个给定的
$[z_0,w_0]\in\operatorname{Crit}\mathcal A_{H^0}$ 与
$[z_1,w_1]\in\operatorname{Crit}\mathcal A_{H^1}$，定义
\[
\begin{aligned}
 &\mathcal M((\mathcal H,\mathfrak j;\rho);[z_0,w_0],[z_1,w_1])\\
 &\qquad:=\bigcup_C
 \mathcal M((\mathcal H,\mathfrak j;\rho);z_0,z_1;C),
\end{aligned}
\]
其中 $C\in\pi_2(z_0,z_1)$ 满足
\begin{equation}
 [z_1,w_1]=[z_1,w_0\#C].
\label{eq:19-chain-map-class-condition}
\end{equation}
这与~\eqref{eq:19-boundary-class-condition} 类似。若 $\rho$-拉长
$(\mathcal H^\rho,\mathfrak j^\rho)$ 在定义 19.3.2 的意义下\emph{\personFloer{}正则}，则称
$(\mathcal H,\mathfrak j;\rho)$ \emph{\personFloer{}正则}。

在定义 19.3.2 的条件下，可像边界映射那样，由矩阵元
$n_{(\mathcal H,\mathfrak j;\rho)}([z_0,w_0],[z_1,w_1])$ 定义次数为 $0$ 的映射
\[
 h_{(\mathcal H,\mathfrak j;\rho)}:CF(H^0)\longrightarrow CF(H^1).
\]
定义 19.3.2 的条件还蕴含 $h_{(\mathcal H,\mathfrak j)}$ 次数为 $0$，且满足
\[
 h_{(\mathcal H,\mathfrak j;\rho)}\circ\partial_{(H^0,J^0)}
 =\partial_{(H^1,J^1)}\circ h_{(\mathcal H,\mathfrak j;\rho)}.
\]
任意两个这样的链映射 $h_{(\mathfrak j_1,\mathcal H^1)}$ 与
$h_{(\mathfrak j_2,\mathcal H^2)}$ 也链同伦；下面给出证明。

\section{链同伦映射的构造}
\label{sec:construction-chain-homotopy-map}

虽然上述同调同构只依赖于端点哈密顿量 $H_\alpha,H_\beta$，相应的链映射却不仅依赖
连接 $H_\alpha$ 与 $H_\beta$ 的同伦
$\mathcal H=\{H(\eta)\}_{0\leq\eta\leq1}$，还依赖同伦
$\mathfrak j=\{J(\eta)\}_{0\leq\eta\leq1}$。固定非退化哈密顿量
$H_\alpha,H_\beta$ 以及连接它们的同伦 $\mathcal H$，再固定相容殆复结构的同伦
$\mathfrak j=\{J(\eta)\}_{0\leq\eta\leq1}$ 和截断函数
$\rho:\mathbb R\to[0,1]$。

回忆，在前几节定义模空间
$\mathcal M(H,J;[z^-,w^-],[z^+,w^+])$ 与\allowbreak
$\mathcal M((\mathcal H,\mathfrak j;\rho);[z^-,w^-],[z^+,w^+])$ 时，施加了同伦条件
\begin{equation}
 [z^+,w^+]=[z^+,w^-\#u];\qquad
 [u]=C\quad\text{于 }\pi_2(z^-,z^+).
\label{eq:19-chain-homotopy-class-condition}
\end{equation}
条件~\eqref{eq:19-chain-homotopy-class-condition} 的一个推论是
\[
 [z^+,w^+]=[z^+,w^-\#u]\quad\text{于 }\Gamma,
\]
但后者比前者弱。换言之，可能存在不止一对互异元素
$C_1,C_2\in\pi_2(z^-,z^+)$，使得
\[
 \mu(z^-,z^+;C_1)=\mu(z^-,z^+;C_2),
 \qquad\omega(C_1)=\omega(C_2).
\]

给定同伦的同伦 $(\overline{\mathcal H},\overline{\mathfrak j})$，其中
$\overline{\mathfrak j}=\{\mathfrak j_\kappa\}$、
$\overline{\mathcal H}=\{\mathcal H_\kappa\}$，还定义由
$\mathcal H_\kappa^\rho$ 组成的拉长
\[
 \overline{\mathcal H}^{\rho}
 =\{\mathcal H_\kappa^\rho\}_{0\leq\kappa\leq1}.
\]
对每个 $0\leq\kappa\leq1$，考虑 (18.5.44) 的参数化版本
\begin{equation}
 \left\{
 \begin{aligned}
 &\frac{\partial u}{\partial\tau}
 +J_\kappa^\rho\left(\frac{\partial u}{\partial t}
 -X_{H_\kappa^\rho}(u)\right)=0,\\
 &u(-\infty)=[z^-,w^-],\qquad u(\infty)=[z^+,w^+],
 \end{aligned}
 \right.
\label{eq:19-parameterized-continuation-equation}
\end{equation}
并定义参数化模空间
\[
\begin{aligned}
 &\mathcal M^{\mathrm{para}}
 ((\overline{\mathcal H},\overline{\mathfrak j};\rho);
 [z^-,w^-],[z^+,w^+])\\
 &\quad=\bigcup_{\kappa\in[0,1]}\{\kappa\}\times
 \mathcal M((\mathcal H_\kappa,\mathfrak j_\kappa;\rho);
 [z^-,w^-],[z^+,w^+]).
\end{aligned}
\]
这里要强调，在考察参数化模空间时，渐近条件不依赖参数
$\kappa\in[0,1]$，而是固定的。

考虑这样的对子 $[z^-,w^-],[z^+,w^+]$：相伴参数化模空间满足
\[
 \dim\mathcal M^{\mathrm{para}}
 \bigl((\overline{\mathcal H},\overline{\mathfrak j};\rho);
 [z^-,w^-],[z^+,w^+]\bigr)=0.
\]
借助其基数给出的矩阵元，定义链同伦映射
\footnote{译者注：原书这里把用于计数的空间写成未参数化的
$\mathcal M((\mathcal H_\kappa,\mathfrak j_\kappa;\rho);\cdot,\cdot)$，其中
$\kappa$ 未量化；但前文刚定义参数化模空间，后文又明确由参数空间贡献附加的一维，
故这里依上下文改为 $\mathcal M^{\mathrm{para}}$。}
\[
 \Upsilon_{\overline{\mathcal H}^{\rho}}:
 CF_*(H_\alpha)\longrightarrow CF_*(H_\beta).
\]
（这里仍像此前一样，在半正性假设下使用基于维数计数的横截性--紧性论证。）
由于参数空间贡献一个附加维数，得到
\[
 \mu_{H_\alpha}([z^-,w^-])-\mu_{H_\beta}([z^+,w^+])=-1,
\]
故映射 $\Upsilon_{\overline{\mathcal H}^{\rho}}$ 的次数为 $1$。

再次利用\personGromov{}--\personFloer{}紧性，可得到
$\overline{\mathcal M}^{\mathrm{para}}$ 的边界结构定理。

\begin{proposition}
\label{prop:boundary-parameterized-chain-homotopy-moduli}
假设 $(\mathcal H_0,\mathfrak j_0)$ 与 $(\mathcal H_1,\mathfrak j_1)$ 在定义 19.3.2
的意义下\personFloer{}正则，且 $(\overline{\mathcal H},\overline{\mathfrak j})$ 在相应的
参数化意义下\personFloer{}正则。设
$[z^-,w^-]\in\operatorname{Crit}\mathcal A_{H_\alpha}$ 与
$[z^+,w^+]\in\operatorname{Crit}\mathcal A_{H_\beta}$ 满足
\[
 \mu_{H_\alpha}([z^-,w^-])-\mu_{H_\beta}([z^+,w^+])=0,
\]
从而相伴的参数化模空间维数为 $1$。则对某个整数 $\ell\geq0$，存在有限多个点
$0<\kappa_1,\ldots,\kappa_\ell<1$，并且
\footnote{译者注：源边界式的固定端因子即
$\{0\}\times\mathcal M(H_\alpha,J_\alpha;\cdot)$ 与
$\mathcal M(H_\beta,J_\beta;\cdot)\times\{1\}$，与译文一致；中间段因子源印
覆盖整个 $\kappa$ 区间的 $\mathcal M^{\mathrm{para}}((\overline{\mathcal H},
\overline{\mathfrak j};\rho);\cdot)$。译文按参数化模空间在 $\kappa=0,1$ 端
退化为端点纤维的理解，把中间段改写为
$\mathcal M((\mathcal H_0,\mathfrak j_0;\rho);\cdot)$ 与
$\mathcal M((\mathcal H_1,\mathfrak j_1;\rho);\cdot)$，系解释性重写；源书字面
则保留 $\mathcal M^{\mathrm{para}}$ 因子。末项又把 $\{\kappa_j\}$ 与
$\mathcal H_{\kappa_i}$ 混用，译文统一为 $i$。}
\[
\begin{aligned}
 &\partial\overline{\mathcal M}^{\mathrm{para}}
 ((\overline{\mathcal H},\overline{\mathfrak j};\rho);
 [z^-,w^-],[z^+,w^+])\\
 &=\bigcup_{[z'^-,w'^-]}
 \bigl(\{0\}\times\mathcal M(H_\alpha,J_\alpha;[z^-,w^-],[z'^-,w'^-])\bigr)\\
 &\qquad\times
 \mathcal M((\mathcal H_0,\mathfrak j_0;\rho);
 [z'^-,w'^-],[z^+,w^+])\\
 &\quad\cup\bigcup_{[z'^+,w'^+]}
 \mathcal M((\mathcal H_1,\mathfrak j_1;\rho);
 [z^-,w^-],[z'^+,w'^+])\\
 &\qquad\times
 \bigl(\{1\}\times\mathcal M(H_\beta,J_\beta;[z'^+,w'^+],[z^+,w^+])\bigr)\\
 &\quad\cup\bigcup_{i=1}^{\ell}\{\kappa_i\}\times
 \mathcal M((\mathcal H_{\kappa_i},\mathfrak j_{\kappa_i};\rho);
 [z^-,w^-],[z^+,w^+]).
\end{aligned}
\]
\end{proposition}

于是该命题立即给出如下恒等式（差一个符号）：
\begin{equation}
 h_{(\mathfrak j_1,\mathcal H_1;\rho_1)}
 -h_{(\mathfrak j_0,\mathcal H_0;\rho_0)}
 =\partial_{(J^1,H^1)}\circ\Upsilon_{\overline{\mathcal H}^{\rho}}
 +\Upsilon_{\overline{\mathcal H}^{\rho}}\circ\partial_{(J^0,H^0)}.
\label{eq:19-chain-homotopy-identity}
\end{equation}
映射 $\Upsilon_{\overline{\mathcal H}^{\rho}}$ 仍依赖于同伦
$\overline{\mathfrak j}$ 与 $\rho$ 的选择，故也记作
\[
 \Upsilon_{\overline{\mathcal H}^{\rho}}
 =\Upsilon_{(\overline{\mathcal H},\overline{\mathfrak j};\rho)}.
\]
式~\eqref{eq:19-chain-homotopy-identity} 特别证明了：连接相同端点的不同同伦
$(\mathfrak j_0,\mathcal H_0;\rho)$ 与 $(\mathfrak j_1,\mathcal H_1;\rho)$ 所给出的
两个链映射链同伦。因此，链映射~\eqref{eq:19-floer-homology-continuation-map} \emph{在同调上}
诱导的映射与同伦 $(\mathcal H,\mathfrak j)$ 或 $\rho$ 的选择无关。

\begin{remark}
\label{rem:vary-elongation-in-chain-homotopy}
上面的讨论没有让拉长函数 $\rho$ 随 $\kappa$ 变化，这只是为了简化记号。最自然的
做法其实是让拉长函数依赖于 $\kappa$。
\end{remark}

\section{\personFloer{}链映射的复合法则}
\label{sec:composition-law-floer-chain-maps}

本节考察\personFloer{}同态
\begin{equation}
 h_{\alpha\beta}:HF_*(H_\alpha)\longrightarrow HF_*(H_\beta)
\label{eq:19-floer-homology-continuation-map}
\end{equation}
的复合法则
\[
 h_{\alpha\gamma}=h_{\beta\gamma}\circ h_{\alpha\beta}.
\]
现在重新考察方程 (18.5.42)。研究\personFloer{}模空间的一个关键分析事实是存在先验能量
上界；下一节将对此作说明。

考虑哈密顿量三元组
\[
 (H_\alpha,H_\beta,H_\gamma)
\]
以及同伦 $\mathcal H_1=\{H_1(s)\}_{0\leq s\leq1}$、
$\mathcal H_2=\{H_2(s)\}_{0\leq s\leq1}$，满足
\[
 H_1(0)=H_\alpha,\quad H_1(1)=H_\beta=H_2(0),\quad H_2(1)=H_\gamma.
\]
定义它们的拼接
$\mathcal H_1\#\mathcal H_2=\{H_3(s)\}_{0\leq s\leq1}$ 为
\[
 H_3(s)=
 \begin{cases}
  H_1(2s),&0\leq s\leq\frac12,\\
  H_2(2s-1),&\frac12\leq s\leq1.
 \end{cases}
\]

注意，由于截断函数 $\rho$ 的选取，延拓方程 (18.5.42) 在区域 $|\tau|>R$ 上是
\emph{自治}的，即在 $\tau$ 的平移下不变。给定三元组
$(H_\alpha,H_\beta,H_\gamma)$ 时，这一事实使我们能够分别粘合对应于
$(H_\alpha,H_\beta)$ 与 $(H_\beta,H_\gamma)$ 的两个这种方程的解。

下面作更精确的说明。给定一对拉长函数
\[
 \rho=(\rho_1,\rho_2)
\]
和正数 $R>0$，定义 $\mathcal H_1\#\mathcal H_2$ 的拉长同伦
\[
 \mathcal H_1\#_{(\rho;R)}\mathcal H_2
 =\{H_{(\rho;R)}(\tau)\}_{-\infty<\tau<\infty}
\]
如下：
\[
 H_{(\rho;R)}(\tau,t,x)=
 \begin{cases}
 H_1(\rho_1(\tau+2R),t,x),&\tau\leq0,\\
 H_2(\rho_2(\tau-2R),t,x),&\tau\geq0.
 \end{cases}
\]
注意，对分别依赖于拉长函数 $\rho_1,\rho_2$ 与同伦
$\mathcal H_1,\mathcal H_2$ 的某些充分大的 $R_1,R_2>0$，有
\[
 H_{(\rho;R)}\equiv
 \begin{cases}
 H_\alpha,&\tau\leq-(R_1+2R),\\
 H_\beta,&-R\leq\tau\leq R,\\
 H_\gamma,&\tau\geq R_2+2R.
 \end{cases}
\]
\emph{特别地，即使通常的粘合同伦 $\mathcal H_1\#\mathcal H_2$ 未必光滑，这个拉长同伦
也总是光滑的}。类似地定义 $\mathfrak j_1\#_{(\rho;R)}\mathfrak j_2$。

对拉长同伦
$(\mathfrak j_1\#_{(\rho;R)}\mathfrak j_2,
\mathcal H_1\#_{(\rho;R)}\mathcal H_2)$，考虑相应的扰动\personCauchy{}--\personRiemann{}方程
\[
 \left\{
 \begin{aligned}
 &\frac{\partial u}{\partial\tau}
 +J_3^{\rho(\tau)}
 \left(\frac{\partial u}{\partial t}-X_{H_3^{\rho(\tau)}}(u)\right)=0,\\
 &\lim_{\tau\to-\infty}u(\tau)=z^-,\qquad
 \lim_{\tau\to\infty}u(\tau)=z^+,
 \end{aligned}
 \right.
\]
并附加条件~\eqref{eq:19-chain-homotopy-class-condition}。

设 $u_1,u_2$ 分别是与 $\rho_1,\rho_2$ 对应的方程 (18.5.44) 的解。在区域
$-R\leq\tau\leq R$ 上用一个单位分解作适当的固定插值，并在两端令预粘合映射为
\[
 u_1\#_Ru_2(\tau,t)=
 \begin{cases}
 u_1(\tau+2R,t),&\tau\leq-R,\\
 u_2(\tau-2R,t),&\tau\geq R.
 \end{cases}
\]
若 $R_0$ 充分大，则赋值
\[
 (u_1,u_2,R)\longmapsto u_1\#_Ru_2
\]
定义从
\[
\begin{aligned}
 &\mathcal M(\mathfrak j_1,\mathcal H_1;[z_1,w_1],[z_2,w_2])
 \times\mathcal M(\mathfrak j_2,\mathcal H_2;[z_2,w_2],[z_3,w_3])\\
 &\hspace{8em}\times(R_0,\infty)
\end{aligned}
\]
到其像的微分同胚。

记 $\overline\partial_{(\mathcal H,\mathfrak j;\rho)}$ 为相应的扰动\personCauchy{}--\personRiemann{}算子
\[
 u\longmapsto\frac{\partial u}{\partial\tau}
 +J_3^{\rho(\tau)}
 \left(\frac{\partial u}{\partial t}-X_{H_3^{\rho(\tau)}}(u)\right),
\]
它作用在满足渐近条件 $u(\pm\infty)=z^\pm$ 和固定同伦条件
$[u]=C\in\pi_2(z^-,z^+)$ 的映射上。在适当的 $u$ 的 $W^{1,p}$ 空间（$p>2$）中，
把 $u_1\#_Ru_2$ 扰动一个小于它成为真解所差误差的量，也就是小于加权 $L^p$-范数
\[
 \bigl\|\overline\partial_{(\mathcal H,\mathfrak j;\rho)}
 (u_1\#_{(\rho;R)}u_2)\bigr\|_p,
\]
便可在 $u_1\#_Ru_2$ 附近构造唯一的真解。利用第 15.5 节的粘合定理，略用记号，
仍把这个真解记作 $u_1\#_Ru_2$。相应映射于是给出嵌入
\[
\begin{aligned}
 &\mathcal M(\mathfrak j_1,\mathcal H_1;[z_1,w_1],[z_2,w_2])
 \times\mathcal M(\mathfrak j_2,\mathcal H_2;[z_2,w_2],[z_3,w_3])
 \times(R_0,\infty)\\
 &\quad\longrightarrow
 \mathcal M(\mathfrak j_1\#_{(\rho;R)}\mathfrak j_2,
 \mathcal H_1\#_{(\rho;R)}\mathcal H_2;[z_1,w_1],[z_3,w_3]).
\end{aligned}
\]
特别地，当
\[
 \mu_{H_\beta}([z_2,w_2])-\mu_{H_\alpha}([z_1,w_1])
 =\mu_{H_\gamma}([z_3,w_3])-\mu_{H_\beta}([z_2,w_2])=0
\]
时，$\mathcal M(\mathfrak j_1,\mathcal H_1;[z_1,w_1],[z_2,w_2])$ 与
$\mathcal M(\mathfrak j_2,\mathcal H_2;[z_2,w_2],[z_3,w_3])$ 都紧致，因而由有限多个点
组成。此外，上述嵌入的像穷尽模空间
\[
 \mathcal M(\mathfrak j_1\#_{(\rho;R)}\mathfrak j_2,
 \mathcal H_1\#_{(\rho;R)}\mathcal H_2;[z_1,w_1],[z_3,w_3])
\]
的“端”，其紧化的边界由断裂流线
\[
 u_1\#_{(\rho;\infty)}u_2=u_1\#_\infty u_2
\]
组成。这就证明了如下粘合恒等式。

\begin{proposition}
\label{prop:composition-gluing-chain-maps}
存在 $R_0>0$
\footnote{译者注：原书此式右端把两个链映射的复合次序写成
$h_{(\mathcal H_1,\mathfrak j_1;\rho_1)}\circ
h_{(\mathcal H_2,\mathfrak j_2;\rho_2)}$，但这与陈述中的定义域、陪域不匹配，
也与本节开头的 $h_{\alpha\gamma}=h_{\beta\gamma}\circ h_{\alpha\beta}$ 相反；
译文按后者改正。}
，使得对任意 $R\geq R_0$，作为从 $CF_*(H_\alpha)$ 到
$CF_*(H_\gamma)$ 的链映射，有
\[
 h_{(\mathcal H_1,\mathfrak j_1)\#_{(\rho;R)}
 (\mathcal H_2,\mathfrak j_2)}
 =h_{(\mathcal H_2,\mathfrak j_2;\rho_2)}
 \circ h_{(\mathcal H_1,\mathfrak j_1;\rho_1)}.
\]
\end{proposition}

这里提醒读者，同伦 $\mathcal H_1\#_{(\rho;R)}\mathcal H_2$ 本身是粘合同伦
$\mathcal H_1\#\mathcal H_2$ 的一个拉长。此命题于是给出同调上的复合法则
$h_{\alpha\gamma}=h_{\beta\gamma}\circ h_{\alpha\beta}$。

\section{横截性}
\label{sec:hamiltonian-floer-transversality}

一般而言，通过通有地选取哈密顿量 $H$ 或殆复结构 $J$，总能使\emph{光滑}\personFloer{}模空间达到横截。
但对紧化\personFloer{}模空间却不能如此，因为会出现负\personChern{}数多重覆叠球面的冒泡。一般
必须在\termKuranishiStructure{}的抽象框架中，利用虚基本链与多值扰动技术取得定义各类\personFloer{}映射所需的一般横截性。

\stfhIndex{I112}{Huofeier-Salameng}{霍费尔--萨拉蒙（Hofer--Salamon）}
我们把讨论限制在半正情形，依照\personHofer{}与\personSalamon{} (HS95) 的叙述，说明半正情形下
如何满足定义 19.2.7 和 19.3.2 中的假设。借助\termKuranishiStructure{}处理一般情形留给
(FOn99)。先给出 (HS95)、(On95) 中半正辛流形的定义。

\begin{definition}
\label{def:semipositive-symplectic-manifold}
若不存在 $A\in\pi_2(M)$ 使得
\[
 \omega(A)>0\qquad\text{且}\qquad 3-n\leq c_1(A)<0,
\]
则称辛流形 $(M,\omega)$ 是\emph{半正}的。
\end{definition}

对边界映射 $\partial_{(H,J)}$，有如下定理。

\begin{theorem}
\label{thm:semipositive-boundary-regularity}
假设不存在 $A\in\pi_2(M)$ 使得
\[
 \omega(A)>0\qquad\text{且}\qquad 3-n\leq c_1(A)<0.
\]
考虑非退化哈密顿量 $H:S^1\times M\to\mathbb R$。则对 $H$-正则的 $J$，
定义 19.2.7 中的假设成立。
\end{theorem}

\begin{proof}
边界映射 $\partial_{(H,J)}$ 的矩阵元由相对\personConley{}--\personZehnder{}指标满足
\begin{equation}
 \mu_H([z^-,w^-],[z^+,w^+])=1
\label{eq:19-index-one-boundary-moduli}
\end{equation}
的模空间
$\widetilde{\mathcal M}(H,J;[z^-,w^-],[z^+,w^+])$ 决定。我们已经证明，该空间是
有限多个 $C\in\pi_2(z^-,z^+)$ 所对应的
$\widetilde{\mathcal M}(H,J;z^-,z^+;C)$ 的并。

目标是证明
$\mathfrak j_\omega=C^\infty([0,1],\mathcal J_\omega)$ 中存在稠密的
$H$-正则 $J$ 子集。先通有地选取 $J$，使
$\widetilde{\mathcal M}(H,J;z^-,z^+;C)$ 在
$\mu_{(H,J)}(C)\leq0$ 时为空，在 $\mu_{(H,J)}(C)>0$ 时具有预期维数。这来自通有横截
性，其证明与第 10.4 节相同。
\footnote{译者注：原书这里笼统声称指标不大于零时模空间为空，但定义 19.2.7(1)
明确保留 $z_0=z_1,C=0$ 的定常指标零解。此句若不排除该情形便与定义冲突；译文
照录，并在此指出缺失的例外。}

\begin{exercise}
\label{ex:prove-hamiltonian-generic-transversality}
证明此横截性。
\end{exercise}

对这样的通有 $J$，$\widetilde{\mathcal M}(H,J;z^-,z^+;C)$ 是光滑一维流形。
指标条件~\eqref{eq:19-index-one-boundary-moduli} 与指标公式 (18.4.25)
\[
 \mu_H([z,w\#A])=\mu_H([z,w])-2c_1(A)
\]
蕴含 $z^-\neq z^+$，所以对每个这样的 $C$，$\mathbb R$ 在
$\widetilde{\mathcal M}(H,J;z^-,z^+;C)$ 上自由作用。故商
\[
 \mathcal M(H,J;z^-,z^+;C)
 =\widetilde{\mathcal M}(H,J;z^-,z^+;C)/\mathbb R
\]
是零维流形。矩阵元
\[
 n_{(H,J)}([z^-,w^-],[z^+,w^+])
\]
应当是这些 $C$ 所对应的基数
\[
 \#\mathcal M(H,J;z^-,z^+;C)
\]
之和。为了定义这个数，需要证明
$\widetilde{\mathcal M}(H,J;z^-,z^+;C)/\mathbb R$ 紧致。

\begin{proposition}
\label{prop:index-one-semipositive-compactness}
假设 $\mu_H([z^-,w^-],[z^+,w^+])=1$。则对通有的 $J$，
$\mathcal M(H,J;[z^-,w^-],[z^+,w^+])$ 紧致。
\end{proposition}

\begin{proof}
由\personGromov{}--\personFloer{}紧性，
$\widetilde{\mathcal M}(H,J;z^-,z^+;C)/\mathbb R$ 中的任意发散子列，一般有子列收敛到
\[
 u_\infty=u_1\#\cdots\#u_N;\qquad
 u_1(-\infty)=z^-,\quad u_N(\infty)=z^+.
\]
这是由尖点\personFloer{}流线 $u_\ell$ 组成的连通链，满足
\[
 u_\ell^0(-\infty)=z_\ell^-,\qquad
 u_\ell^0(\infty)=z_\ell^+;\qquad
 z_\ell^+=z_{\ell+1}^-,
\]
其中 $1\leq\ell\leq N$，并且每个 $u_\ell$ 形如
\begin{equation}
 u_\ell=u_\ell^0+\sum_{k=1}^{L_\ell}v_\ell^k.
\label{eq:19-cusp-floer-components}
\end{equation}
这里 $\{v_\ell^k\}_k$ 是有限个球面气泡，记
$C_\ell=[u_\ell]\in\pi_2(z_\ell^-,z_\ell^+)$。
\footnote{译者注：原书式 (19.6.21) 对每个分支都以上限 $L$ 计数气泡，但后文立即
改用逐分支的 $L_j$ 并令 $L=\sum_jL_j$；译文统一把本式上限写作 $L_\ell$。}

为证明 $\mathcal M(H,J;z^-,z^+;C)$ 紧致，我们要证明：对通有的 $J$，当
$\mu_{(H,J)}(C)=1$ 时，对任意 $(z^-,z^+)$ 与
$C\in\pi_2(z^-,z^+)$，都有 $N=1$ 且 $\{v_\ell^k\}=\varnothing$。

回忆，每个 $v_\ell^k$ 是在某点 $(\tau_k,t_k)$ 附着于 $u_\ell$ 的
$J_{t_k}$-全纯球面，其中先验地也允许 $\tau_k=\infty$。先考察 $\tau_k=\infty$，
即气泡附着在周期轨道像集上的情形。

考虑参数化模空间
\[
 \mathcal M^{\mathrm{para}}(J;\alpha)
 =\bigcup_{t\in[0,1]}\{t\}\times\mathcal M(J_t;\alpha).
\]
它的虚维数为
\[
 2c_1(\alpha)+2n-6+1=2c_1(\alpha)+2n-5.
\]
对通有 $J$，$\mathcal M^{\mathrm{para}}(J;\alpha)$ 中在某处单射曲线组成的开子集是
维数 $2c_1(\alpha)+2n-5$ 的光滑流形。因此，只要该模空间非空，就必有
\begin{equation}
 c_1(\alpha)\geq3-n.
\label{eq:19-simple-sphere-chern-lower-bound}
\end{equation}
现在假设 $J$ 对所有在某处单射曲线都满足此性质。

另一方面，对多重覆叠气泡，特别是\personChern{}数满足
$3-n\leq c_1(\alpha)<0$ 的气泡，一般不能通过扰动 $J$ 获得所需横截性。\emph{这里半正性
起关键作用}。在半正性假设下，式~\eqref{eq:19-simple-sphere-chern-lower-bound} 自动蕴含
$c_1(\alpha)\geq0$。若通有地选取 $J$，使该式对所有在某处单射曲线成立，那么由结构
定理 8.7.3，任何不在某处单射的曲线 $u$ 都是多重覆叠：
$u=\widetilde u\circ\varphi$，其中 $\varphi:S^2\to S^2$ 是次数 $d>1$ 的分歧覆叠，
$\widetilde u$ 在某处单射。因此
\[
 c_1(u)=d\,c_1(\widetilde u)\geq c_1(\widetilde u)
\]
（等号仅在 $c_1(\widetilde u)=0$ 时成立），从而
\[
 \operatorname{vir.dim}\mathcal M(J;[u])
 \geq\operatorname{vir.dim}\mathcal M(J;[\widetilde u]).
\]
显然 $\widetilde u$ 与 $u$ 的像集相同，所以若 $\widetilde u$ 不与
$\operatorname{Per}(H)$ 相交，$u$ 的像也不与之相交。这里略用记号，
$\operatorname{Per}(H)\subset M$ 表示 $H$ 的周期轨道的像集。因此，只需关注在某处
单射的气泡。

现在考虑评价映射
\[
 \operatorname{ev}_0:\mathcal M_1^{\mathrm{para}}(J;\alpha)
 \longrightarrow M,\qquad \operatorname{ev}_0(v,z)=v(z).
\]
$\mathcal M_1^{\mathrm{para}}(J;\alpha)$ 的虚维数比
$\mathcal M^{\mathrm{para}}(J;\alpha)$ 大 $2$，故为
$2c_1(\alpha)+2n-3$。必要时进一步扰动 $J$ 后，评价映射在在某处单射曲线的开子集上
与给定的一维子流形 $\operatorname{Per}(H)\subset M$ 横截。因此若
$3-2c_1(\alpha)>1$，即若
\begin{equation}
 c_1(\alpha)<1,
\label{eq:19-evaluation-avoids-periodic-orbits}
\end{equation}
则 $\operatorname{ev}_0$ 的像不与 $\operatorname{Per}(H)$ 相交。更多气泡更难出现的
理由留给读者说明。

结合~\eqref{eq:19-simple-sphere-chern-lower-bound} 与
\eqref{eq:19-evaluation-avoids-periodic-orbits} 可知，若
$3-n\leq c_1(\alpha)\leq0$，则 $X_H$ 的任意周期轨道的像上都不会发生冒泡。

由于每个 $u_k$ 的指标应至少为 $1$，有 $\mu_{(H,J)}([u_k])\geq1$。既然不能在
$(\pm\infty,*)$ 处冒泡，冒泡只能发生在某点
$(\tau_\ell^k,t_\ell^k)\in\mathbb R\times\mathbb R/\mathbb Z$。于是考虑评价映射
\[
 \operatorname{ev}_C:\mathcal M_L(H,J;z^-,z^+;C)\longrightarrow M^L,
\]
其中 $L$ 是标记点个数，$\mathcal M_L(H,J;z^-,z^+;C)$ 是带 $L$ 个标记点的\personFloer{}流线模空间。更确切地，
\[
 \mathcal M_L(H,J;z^-,z^+;C)
 =\widetilde{\mathcal M}(H,J;z^-,z^+;C)
 \times(\mathbb R\times S^1)^L/\sim,
\]
其中等价关系 $\sim$ 定义为
\footnote{译者注：原书定义等价关系时把标记点列写成
$(z_1,\ldots,z_L)$，下一行却立即以坐标 $(\tau_i,t_i)$ 使用它们；译文统一写出这些
标记点的柱面坐标。}
\[
 (u;((\tau_1,t_1),\ldots,(\tau_L,t_L)))
 \sim(u';((\tau'_1,t'_1),\ldots,(\tau'_L,t'_L)))
\]
当且仅当对某个 $\tau^0$，
\[
 u'(\tau,t)=u(\tau-\tau^0,t),\qquad
 (\tau_i',t_i')=(\tau_i+\tau^0,t_i).
\]

由 $\sim$ 的定义，$\operatorname{ev}_C$ 良定义。式~\eqref{eq:19-cusp-floer-components}
中的每个位形都是纤维积
\[
 \mathcal M_{L_j}(H,J;z_j,z_{j+1};C_j)_{\operatorname{ev}_{C_j}}
 \!\times_{(\operatorname{ev}_{\alpha_1},\ldots,
 \operatorname{ev}_{\alpha_{L_j}})}
 \bigl(\mathcal M_1^{\mathrm{para}}(\alpha_1^j),\ldots,
 \mathcal M_1^{\mathrm{para}}(\alpha_{L_j}^j)\bigr)
\]
中的元素，并且
\footnote{译者注：原书此处写
$C=C^0\#\sum_{j=1}^N\sum_{k=1}^{L_j}\alpha_k^j$，但本证明未定义 $C^0$，而紧接着
的指标公式逐一使用 $C_j$。译文依极限的 $N$ 段分解显式写出
$C_1\#\cdots\#C_N$。}
\[
 C=C_1\#\cdots\#C_N\#
 \sum_{j=1}^{N}\sum_{k=1}^{L_j}\alpha_k^j.
\]
当 $N\geq1$ 时，这个纤维积的虚维数为
\begin{equation}
 \sum_{j=1}^{N}\bigl(\mu_H(C_j)+2L_j\bigr)-(N-1)
 +\sum_{j=1}^{N}\sum_{k=1}^{L_j}
 \bigl(2c_1(\alpha_k^j)+2n-3\bigr)-2nL,
\label{eq:19-bubbled-stratum-virtual-dimension}
\end{equation}
其中 $L_j\geq0$ 且 $L:=\sum_{j=1}^{N}L_j$。（若 $N=0$，则必有 $z^-=z^+$；
由 $J$ 的选取，不可能附着气泡。此时，对某个 $A\in\pi_2(M)$ 必有
\[
 [z^+,w^+]=[z^-,w^-\#A],\qquad C=C^0=-A.
\]
于是由指标公式，$\mu_H(z^-,z^+;C^0)=2c_1(A)$，不可能等于所要求的 $1$。）
此外，
\begin{equation}
 \mu_H(C)=\sum_{j=1}^{N}
 \left(\mu_H(C_j)+\sum_{k=1}^{L_j}2c_1(\alpha_k^j)\right).
\label{eq:19-total-index-bubble-decomposition}
\end{equation}

\begin{lemma}
\label{lem:index-one-bubbled-stratum-dimension}
假设 $\mu_H(C)=1$。则~\eqref{eq:19-bubbled-stratum-virtual-dimension} 中的虚维数为
$2-L-N$，其中 $N\geq1$。
\end{lemma}

\begin{proof}
把 $\mu_H(C)=1$ 代入~\eqref{eq:19-total-index-bubble-decomposition}，得到
\[
 \sum_{j=1}^{N}
 \left(\mu_H(C_j)+\sum_{k=1}^{L_j}2c_1(\alpha_k^j)\right)=1.
\]
再代入~\eqref{eq:19-bubbled-stratum-virtual-dimension}，并使用
$L=\sum_{j=1}^{N}L_j$，虚维数化为
\[
 1+2L-(N-1)+L(2n-3)-2nL=2-L-N.
\]
\end{proof}

若模空间横截，这个维数必须非负，故 $0\leq2-N-L$。因为 $N\geq1$，$L$ 只能为
$0$ 或 $1$。当 $L=0$ 时，$1=\sum_{j=1}^{N}\mu_H(C_j)$。又因
$\mu_H(C_j)\geq0$，除一个 $j$（不妨为 $j=1$）外，其余 $\mu_H(C_j)$ 都为 $0$。
于是对所有 $j\neq1$ 都有 $z_j^-=z_j^+$，且对 $j=2,\ldots,N$，至少有一个气泡
必须附着于 $z_j^-=z_j^+$；由 $J$ 的选取，这不可能。因此此时也有 $N=1$。

若 $L=1$，上述不等式给出 $N\leq1$，故必有 $N=1$。这意味着气泡附着于一个周期
轨道，但这已经被前面对 $J$ 的选取排除。因此必有 $N=1$、$L=0$；也就是说，
该序列既不产生非常值气泡，也不分裂为断裂流线。这就证明了当
$\mu_H(z^-,z^+;C)=1$ 时，每个 $\mathcal M(H,J;z^-,z^+;C)$ 都紧致。
\footnote{译者注：本段的 $L$ 是有限柱面点处的标记点（冒泡点）总数，所以
$N=1,L=1$ 描述的气泡一般附着于\personFloer{}流线内部，并不自动附着于周期轨道。
原书从 $N=1,L=1$ 直接推出“气泡附着于周期轨道”，与此前的标记点模空间设置
不符；正文照录该推断，但紧性证明在此存在缺口。}

由此命题，边界映射 $\partial_{(H,J)}$ 的矩阵元良定义，因而边界映射本身良定义。
类似的论证证明，当 $\mu_H(z^-,z^+;C)=2$ 时，
$\mathcal M(H,J;z^-,z^+;C)$ 除了上述讨论中 $N=2,L=0$ 的分裂外是紧致的。这就
证明了定义 19.2.7 中的假设，从而定理 19.6.2 得证。
\end{proof}
\end{proof}

\begin{exercise}
\label{ex:index-two-boundary-compactification-details}
补全上述证明最后一段中关于 $\mu_H(z^-,z^+;C)=2$ 情形的论证细节。
\end{exercise}

下面考虑链映射 $h_{(\mathcal H,\mathfrak j;\rho)}$。唯一的差别是这里不取
$\mathbb R$ 的商，而直接考虑模空间
\[
 \mathcal M((\mathcal H,\mathfrak j;\rho);[z_0,w_0],[z_1,w_1])
 =\bigcup_C\mathcal M((\mathcal H,\mathfrak j;\rho);z_0,z_1;C).
\]
其维数为 $\mu_{(\mathcal H,\mathfrak j;\rho)}(C)$，而定义链映射所需的是
$\mu_{(\mathcal H,\mathfrak j;\rho)}(C)=0$ 的情形。类似地可以证明如下定理。

\begin{theorem}
\label{thm:semipositive-chain-map-regularity}
假设 $(M,\omega)$ 中不存在 $A\in\pi_2(M)$ 使得
\[
 \omega(A)>0\qquad\text{且}\qquad3-n\leq c_1(A)<0.
\]
则定义 19.3.2 中的假设成立。
\end{theorem}

\begin{proof}
由同类论证，对 $N\geq1$ 的极限位形得到虚维数
\[
 \sum_{j=1}^{N}(\mu_H(C_j)+2L_j)-(N-1)
 +\sum_{j=1}^{N}\sum_{k=1}^{L_j}
 (2c_1(\alpha_k^j)+2n-3)-2nL,
\]
以及
\[
 \mu_H(C)=\sum_{j=1}^{N}\mu_H(C_j)
 +\sum_{j=1}^{N}\sum_{k=1}^{L_j}2c_1(\alpha_k^j).
\]
若 $N=0$，则必有
\[
 [z^-,w^-]=[z^+,w^+],\qquad C^0=0,
\]
因为极限位形中不可能附着任何气泡。因此，相应模空间
$\mathcal M_{(\mathcal H,\mathfrak j;\rho)}([z^-,w^-],[z^+,w^+])$
\footnote{译者注：链映射连接的两端一般分别属于
$\operatorname{Crit}\mathcal A_{H^0}$ 与
$\operatorname{Crit}\mathcal A_{H^1}$，原书却在 $N=0$ 时直接写
$[z^-,w^-]=[z^+,w^+]$，未说明如何识别这两个不同临界集；正文照录。}
紧致。

现设 $N\geq1$。这次考察 $\mu_H(C)=0$，得到
\[
 0=1-N-L.
\]
这蕴含 $L=0$ 且 $N=1$，进而
$\mathcal M_{(\mathcal H,\mathfrak j;\rho)}([z^-,w^-],[z^+,w^+])$ 紧致。
\end{proof}

接着研究 $h_{(\mathcal H,\mathfrak j;\rho)}$ 的链映射性质，即恒等式
\begin{equation}
 h_{(\mathcal H,\mathfrak j;\rho)}\circ\partial_{(H^0,J^0)}
 =\partial_{(H^1,J^1)}\circ h_{(\mathcal H,\mathfrak j;\rho)}.
\label{eq:19-chain-map-property}
\end{equation}
为此，考虑使链映射模空间
$\mathcal M_{(\mathcal H,\mathfrak j;\rho)}([z^-,w^-],[z^+,w^+])$ 虚维数为 $1$
的对子 $[z^-,w^-],[z^+,w^+]$，并分析其紧化边界。

对链同伦映射，考虑参数化模空间
\[
 \mathcal M^{\mathrm{para}}
 (\overline{\mathcal H},\overline{\mathfrak j};\rho;
 [z^-,w^-],[z^+,w^+]),
\]
其中
$\overline{\mathcal H}=\{\mathcal H_\kappa\}_{0\leq\kappa\leq1}$ 与
$\overline{\mathfrak j}=\{\mathfrak j_\kappa\}_{0\leq\kappa\leq1}$ 是单参数族，
\[
 \mathcal H_\kappa=\{H^{s,\kappa}\},\qquad
 \mathfrak j_\kappa=\{J^{s,\kappa}\};\qquad
 H^{0,\kappa}=H^\alpha,\quad H^{1,\kappa}=H^\beta
\]
对所有 $0\leq\kappa\leq1$ 成立。考虑方程族
\[
 \frac{\partial u}{\partial\tau}
 +J_t^{\rho(\tau),\kappa}
 \left(\frac{\partial u}{\partial t}-X_{H_t^{\rho(\tau),\kappa}}\right)=0,
 \qquad0\leq\kappa\leq1.
\]
考察这样的对子 $([z^-,w^-],[z^+,w^+])$：相应参数化模空间的虚维数为 $0$ 或
$1$，即
\[
 \mu_{H^\alpha}([z^-,w^-])-\mu_{H^\beta}([z^+,w^+])+1=0
 \quad\text{或 }1.
\]
维数为 $0$ 时，
\[
 \mu_{H^\alpha}([z^-,w^-])-\mu_{H^\beta}([z^+,w^+])=-1.
\]

\begin{lemma}
\label{lem:parameterized-index-minus-one-compactness}
对通有族 $(\overline{\mathcal H},\overline{\mathfrak j})$，当
\[
 \mu_{H^\alpha}([z^-,w^-])-\mu_{H^\beta}([z^+,w^+])=-1
\]
时，$\mathcal M^{\mathrm{para}}
(\overline{\mathcal H},\overline{\mathfrak j};\rho;
[z^-,w^-],[z^+,w^+])$ 紧致。
\end{lemma}

\begin{exercise}
\label{ex:fredholm-setup-parameterized-compactness}
表述正确的\personFredholm{}设置并证明此引理。
\end{exercise}

借助与这一同伦相伴的族 $\overline{\mathfrak j}$ 与
$\overline{\mathcal H}$，通过计数满足上式指标条件的参数化模空间中的元素，构造链同伦
映射
\[
 \Upsilon_{(\overline{\mathfrak j},\overline{\mathcal H};\rho)}:
 CF_*(J^0,H^0)\longrightarrow CF_{*+1}(J^1,H^1).
\]
特别地，映射 $\Upsilon_{(\overline{\mathfrak j},\overline{\mathcal H};\rho)}$ 的次数为
$1$。
\footnote{译者注：原书同一行把映射的陪域写为 $CF_{*+1}$，且此前第 19.4 节也由
指标差 $-1$ 推出链同伦映射次数为 $1$；但紧接着却称其“次数为 $-1$”。译文按映射
类型和第 19.4 节改为 $1$。}

现在对满足
\[
 \mu_{H^\alpha}([z^-,w^-])-\mu_{H^\beta}([z^+,w^+])=0
\]
的对子考虑上述参数化模空间。此时它的维数为 $1$，其边界如命题 19.4.1 所述。

\begin{remark}
\label{rem:general-case-virtual-moduli-cycle}
对一般辛流形，需要在\termKuranishiStructure{}的框架中使用虚模闭链和抽象多值扰动的概念
(FOn99)。
\end{remark}

\section{时间反演流与对偶性}
\label{sec:time-reversal-flow-duality}

本节研究\personFloer{}复形中的两种对偶配对：一种由流的时间反演诱导，不使用环境流形的
定向；另一种由\personPoincare{}对偶诱导，本质上使用环境流形的定向。

\subsection{时间反演对偶}
\label{subsec:time-reversal-duality}

本小节通过明确考察对\personMorse{}函数及定向系统的依赖性，比较经典\personMorse{}复形上的两种
不同相容定向系统。这里不使用 $M$ 的可定向性。

设
\[
 \sigma(f)=\{o_\sigma(x;f);o_\sigma(x,y;f)\}_{x,y\in\operatorname{Crit}f}
\]
是第 15.6.1 节定义的 $f$ 的\personMorse{}复形上的相容定向系统：\emph{先定向不稳定流形，
再给不稳定流形配上余定向}。它诱导\personMorse{}边界映射
$\partial_f^\sigma:CM_k(f)\to CM_{k-1}(f)$。以
\[
 \delta_f:CM^*(f)\longrightarrow CM^{*+1}(f)
\]
表示对偶复形的相应上边界映射，其中
\[
 CM^\ell(f)=\operatorname{Hom}(CM_\ell(f),\mathbb Z).
\]

回忆，作为集合有 $\operatorname{Crit}f=\operatorname{Crit}(-f)$，而 $f$ 与 $-f$
的稳定、不稳定流形互换。因此，对每个
$y\in\operatorname{Crit}(-f)$，在
\[
 T_yW_f^u(y)=T_yW_{-f}^s(y)
\]
上赋予定向 $\sigma$。注意，作为集合有 $W_{-f}^s(y)=W_f^u(y)$。

接着，以 $\widetilde\sigma(g)$ 表示如下定向系统：\emph{先定向 $g$ 的稳定流形，再为不稳定
流形配上余定向}。

\begin{lemma}
\label{lem:time-reversed-orientation-coherent}
若 $\sigma(f)$ 相容，则定向系统
\[
 \widetilde\sigma(-f)
 =\{o_{\widetilde\sigma}(y;-f)\}_{y\in\operatorname{Crit}(-f)}
\]
相容。
\end{lemma}

\begin{exercise}
\label{ex:prove-time-reversed-orientation-coherent}
证明此引理。
\end{exercise}

利用这一系统，对每对 $y,x\in\operatorname{Crit}(-f)$，也在
$\mathcal M_g(y,x;-f)$ 上诱导一个定向，进而得到边界映射
\[
 \partial_{-f}^{\widetilde\sigma}:CM_\ell(-f)\longrightarrow CM_{\ell-1}(-f).
\]
按定义，“反演流”给出典范一一对应
\[
 \mathcal M_g(x,y;f)\longleftrightarrow\mathcal M_g(y,x;-f).
\]
特别地，$\dim\mathcal M_g(x,y;f)=0$ 当且仅当
$\dim\mathcal M_g(y,x;-f)=0$。此时可以比较两个整数
$n^\sigma(x,y;f)$ 与 $n^{\widetilde\sigma}(y,x;-f)$。

回忆有典范同构
\[
 CM^\ell(f)\longrightarrow CM_{n-\ell}(-f),
\]
它是基元素赋值 $x^*\mapsto\widetilde x$ 的线性延拓；这里
$x\in\operatorname{Crit}f$、$\mu_f(x)=\ell$，$x^*$ 是其对偶，而
$\widetilde x$ 是把 $x$ 看作 $-f$ 的临界点。

\begin{proposition}
\label{prop:time-reversal-morse-sign}
假设 $\mu_f(x)-\mu_f(y)=1$ 且 $\mu_f(y)=k$，则
\[
 n^{\widetilde\sigma}(y,x;-f)=(-1)^k n^\sigma(x,y;f).
\]
特别地，有交换图
\begin{equation}
\begin{array}{ccc}
 CM^k(f)&\xrightarrow{\;(-1)^k\delta_f\;}&CM^{k+1}(f)\\[4pt]
 \big\downarrow&&\big\downarrow\\[4pt]
 CM_{n-k}(-f)&\xrightarrow{\;\partial_{-f}^{\widetilde\sigma}\;}&
 CM_{n-k-1}(-f).
\end{array}
\label{eq:19-time-reversal-commutative-diagram}
\end{equation}
\end{proposition}

\begin{proof}
令 $p\in S W_{-f}^u(y)\cap S W_{-f}^s(x)$。点 $p$ 关于相容定向
$\widetilde\sigma$ 的符号对应于 $o_{\widetilde\sigma}(y,x;-f)$，并由定向向量空间等式
\begin{equation}
 T_p^{\widetilde\sigma}(S W_{-f}^u(y))
 =o_{\widetilde\sigma}(y,x;-f)\,
 N_p^{\widetilde\sigma}(S W_{-f}^s(x))
\label{eq:19-time-reversal-sign-definition}
\end{equation}
决定。这里参见式 (15.6.46) 附近的讨论。
\footnote{译者注：原书在本证明的定向符号中反复写
$o_{\widetilde\sigma}(x,y;-f)$，但相应反演流线是从 $y$ 到 $x$，命题陈述及前文计数
均写 $n^{\widetilde\sigma}(y,x;-f)$；译文统一为 $o_{\widetilde\sigma}(y,x;-f)$。}

另一方面，按 $o_\sigma(x,y;f)$ 的定义，
\begin{equation}
 N_p^\sigma(S W_f^s(y))
 =o_\sigma(x,y;f)\,T_p^\sigma(S W_f^u(x)).
\label{eq:19-original-flow-sign-definition}
\end{equation}
现在比较~\eqref{eq:19-time-reversal-sign-definition} 与
\eqref{eq:19-original-flow-sign-definition}。

式~\eqref{eq:19-time-reversal-sign-definition} 等价于
\begin{equation}
\begin{aligned}
 \{\nabla f(p)\}\oplus N_p^{\widetilde\sigma}(W_{-f}^s(x))
 &=o_{\widetilde\sigma}(y,x;-f)
 \{\nabla f(p)\}\oplus T_p^{\widetilde\sigma}(S W_{-f}^u(y))\\
 &=o_{\widetilde\sigma}(y,x;-f)
 T_p^{\widetilde\sigma}W_{-f}^u(y).
\end{aligned}
\label{eq:19-time-reversal-expanded-orientation}
\end{equation}

另一方面，有如下引理。

\begin{lemma}
\label{lem:stable-minus-f-unstable-f-orientation}
\emph{作为定向向量空间}，
\begin{equation}
 T_p^{\widetilde\sigma}W_{-f}^s(x)=T_p^\sigma W_f^u(x).
\label{eq:19-stable-minus-f-unstable-f}
\end{equation}
\end{lemma}

\begin{proof}
由 $W_{-f}^s(x)$ 的余定向定义，作为定向向量空间有
\[
 N_p^{\widetilde\sigma}W_{-f}^s(x)
 \oplus T_p^{\widetilde\sigma}W_{-f}^s(x)=T_pM.
\]
另一方面，由 $\widetilde\sigma$ 的定义，
\[
 N_p^\sigma W_f^u(x)=N_p^{\widetilde\sigma}W_{-f}^s(x),
\]
故
\[
 N_p^{\widetilde\sigma}W_{-f}^s(x)\oplus T_p^\sigma W_f^u(x)=T_pM.
\]
比较两式即得结论。
\end{proof}

在式~\eqref{eq:19-time-reversal-expanded-orientation} 右侧直和
$T_p^{\widetilde\sigma}(S W_{-f}^s(x))$，得到
\[
\begin{aligned}
 (-1)^{k+1}T_pM
 &=o_{\widetilde\sigma}(y,x;-f)
 T_p^{\widetilde\sigma}W_{-f}^u(y)
 \oplus T_p^{\widetilde\sigma}(S W_{-f}^s(x))\\
 &=o_{\widetilde\sigma}(y,x;-f)
 \{\nabla f(p)\}\oplus T_p^{\widetilde\sigma}(S W_{-f}^u(y))
 \oplus T_p^{\widetilde\sigma}(S W_{-f}^s(x))\\
 &=o_{\widetilde\sigma}(y,x;-f)(-1)^{n-k-1}
 T_p^{\widetilde\sigma}(S W_{-f}^u(y))
 \oplus T_p^{\widetilde\sigma}W_{-f}^s(x).
\end{aligned}
\]
这里左端使用恒等式
\[
\begin{aligned}
 &o_{\widetilde\sigma}(y,x;-f)
 T_p^{\widetilde\sigma}W_{-f}^u(y)
 \oplus T_p^{\widetilde\sigma}(S W_{-f}^s(x))\\
 &\quad=(-1)^{k+1}T_p^{\widetilde\sigma}W_{-f}^u(y)
 \oplus o_{\widetilde\sigma}(y,x;-f)
 \oplus T_p^{\widetilde\sigma}(S W_{-f}^s(x))\\
 &\quad=(-1)^{k+1}T_pM,
\end{aligned}
\]
右端也类似。

在式~\eqref{eq:19-original-flow-sign-definition} 右侧直和
$T_p^\sigma W_f^s(y)$，并注意
$N_p^\sigma(S W_f^s(y))\cong N_p^\sigma W_f^s(y)$，得到
\[
\begin{aligned}
 T_pM
 &=o_\sigma(x,y;f)T_p^\sigma(S W_f^u(x))\oplus T_p^\sigma W_f^s(y)\\
 &=o_\sigma(x,y;f)(-1)^kT_p^\sigma W_f^u(x)
 \oplus T_p^\sigma(S W_f^s(y))\\
 &=o_\sigma(x,y;f)(-1)^k(-1)^{(k+1)(n-k-1)}
 T_p^\sigma(S W_f^s(y))\oplus T_p^\sigma W_f^u(x).
\end{aligned}
\]
故由~\eqref{eq:19-stable-minus-f-unstable-f}，
\[
\begin{aligned}
 &(-1)^{k+1+n-k-1}o_{\widetilde\sigma}(y,x;-f)
 T_p^{\widetilde\sigma}(S W_{-f}^u(y))\\
 &\qquad=(-1)^{k+(k+1)(n-k-1)}o_\sigma(x,y;f)
 T_p^\sigma(S W_f^s(y)).
\end{aligned}
\]
最后注意
\[
 \{\nabla f(p)\}\oplus T_p^{\widetilde\sigma}(S W_{-f}^u(y))
 =T_p^{\widetilde\sigma}W_{-f}^u(y)
\]
以及
\[
 \{-\nabla f(p)\}\oplus T_p^\sigma(S W_f^s(y))=T_p^\sigma W_f^s(y).
\]
把上述关系移到 $y$，得到
\begin{equation}
\begin{aligned}
 &(-1)^{k+1+n-k-1}o_{\widetilde\sigma}(y,x;-f)
 T_y^{\widetilde\sigma}W_{-f}^u(y)\\
 &\qquad=(-1)^{k+(k+1)(n-k-1)+1}o_\sigma(x,y;f)
 T_y^\sigma W_f^s(y).
\end{aligned}
\label{eq:19-orientation-transfer-to-critical-point}
\end{equation}
又有
\[
 N_y^{\widetilde\sigma}W_{-f}^u(y)
 \oplus T_y^{\widetilde\sigma}W_{-f}^u(y)=T_yM
\]
与
\[
 N_y^\sigma W_f^s(y)\oplus T_y^\sigma W_f^s(y)=T_yM.
\]
因此只需比较 $N_y^{\widetilde\sigma}W_{-f}^u(y)$ 与 $N_y^\sigma W_f^s(y)$。
为此直和 $T_y^{\widetilde\sigma}W_{-f}^u(y)$ 并计算
\footnote{译者注：原书这段计算全在临界点 $y$ 的切空间中进行，末项却印作
$T_pM$；译文改为 $T_yM$。}
\[
\begin{aligned}
 N_y^\sigma W_f^s(y)\oplus T_y^{\widetilde\sigma}W_{-f}^u(y)
 &=N_y^\sigma W_f^s(y)\oplus N_y^{\widetilde\sigma}W_{-f}^s(y)\\
 &=T_y^\sigma W_f^u(y)\oplus N_y^\sigma W_f^u(y)\\
 &=(-1)^{k(n-k)}T_yM,
\end{aligned}
\]
其中倒数第二个等号来自 $\widetilde\sigma$ 的定义。这证明
\[
 T_y^{\widetilde\sigma}W_{-f}^u(y)
 =(-1)^{k(n-k)}T_y^\sigma W_f^s(y).
\]
代入~\eqref{eq:19-orientation-transfer-to-critical-point}，得到
\[
 (-1)^{k+1+n-k-1}o_{\widetilde\sigma}(y,x;-f)
 =(-1)^{k+(k+1)(n-k-1)+1+k(n-k)}o_\sigma(x,y;f).
\]
化简指数，便有
\[
 o_{\widetilde\sigma}(y,x;-f)=(-1)^ko_\sigma(x,y;f),
\]
证明完毕。
\end{proof}

由此构造出如下典范同构。

\begin{corollary}[时间反演]
\label{cor:morse-time-reversal-duality}
对任意不必可定向的流形 $M$ 及所有 $k\in\mathbb N$，
\begin{equation}
 HM_{n-k}(-f;\widetilde\sigma)\cong HM^k(f;\sigma).
\label{eq:19-morse-time-reversal-duality}
\end{equation}
\end{corollary}

\subsection{\personMorse{}同调中的\personPoincare{}对偶}
\label{subsec:poincare-duality-morse-homology}

现设 $M$ 可定向，并给定定向 $o_M$。回忆，对有向子流形 $S\subset M$，其法丛在
$p\in S$ 处的定向由定向向量空间关系
\begin{equation}
 N_pS\oplus T_pS=T_pM
\label{eq:19-normal-tangent-ambient-orientation}
\end{equation}
确定。特别地，若按照第 15.6.1 节的方案给 $T_yW_y^u(f)$ 赋予定向，则关系
\[
 N_yW_y^u(f)\oplus T_y^\sigma W_y^u(f)=T_yM
\]
典范地给法空间 $N_yW_y^u(f)$ 本身赋予定向；\emph{没有环境定向时则不能如此}。在当前
情形，由于 $N_yW_f^s(y)=T_yW_f^u(y)$ 已按 $\sigma$ 的定向方案定向，\emph{稳定}流形
$W_f^s(y)$ 也由标准关系
\begin{equation}
 N_y^\sigma W_f^s(y)\oplus T_yW_f^s(y)=T_yM
\label{eq:19-orient-stable-manifold}
\end{equation}
获得自然定向。

把相应的定向向量空间 $T_yW_f^s(y)$ 记作 $T_y^{\check\sigma}W_f^s(y)$。再要求
\begin{equation}
 N^{\check\sigma}W_f^u(y)=T_y^{\check\sigma}W_f^s(y),
\label{eq:19-coorient-unstable-by-stable}
\end{equation}
从而给\emph{不稳定流形}配上余定向。因此，若 $M$ 可定向，则关系
\eqref{eq:19-orient-stable-manifold}--\eqref{eq:19-coorient-unstable-by-stable}
还让 $\sigma(f)$ 产生另一个系统 $\check\sigma(f)$；它给 $f$ 本身而非 $-f$ 的稳定
流形赋予定向，并给不稳定流形赋予余定向。现在比较给定\personMorse{}函数 $f$ 上的
$\sigma(f)$ 与 $\check\sigma(f)$。

系统 $\widetilde\sigma$ 给 $f$ 的稳定流形（或 $-f$ 的不稳定流形）赋予定向。要求
定向向量空间关系
\[
 N_y^{\widetilde\sigma}W_{-f}^u(y)=T_y^{\check\sigma}W_{-f}^s(y),
\]
便在每个 $y$ 处诱导 $N_y^{\check\sigma}W_{-f}^u(y)$ 的一个定向。

\begin{proposition}
\label{prop:compare-two-morse-orientation-systems}
设 $M$ 可定向且 $n=\dim M$，则
\[
 o_{\check\sigma}(x,y;f)=(-1)^no_\sigma(x,y;f).
\]
特别地，符号只依赖于 $\dim M$，而不依赖于同一指标的各个临界点。
\end{proposition}

\begin{proof}
设 $\sigma$ 为上述相容系统。由可定向性和 $\check\sigma$ 的定义，
\[
 T_x^{\check\sigma}W_f^s(x)\oplus T_x^\sigma W_f^u(x)=T_xM,
 \qquad T_x^{\check\sigma}W_f^s(x)=N_x^{\check\sigma}W_f^u(x).
\]
符号 $o_{\check\sigma}(x,y;f)$ 由关系
\[
 N_p^{\check\sigma}(S W_f^u(x))
 =o_{\check\sigma}(x,y;f)T_p^{\check\sigma}(S W_f^s(y))
\]
确定。右侧直和 $T_p^{\check\sigma}(S W_f^u(x))$，得到
\begin{equation}
 T_p^{\check\sigma}M^c
 =o_{\check\sigma}(x,y;f)T_p^{\check\sigma}(S W_f^s(y))
 \oplus T_p^{\check\sigma}(S W_f^u(x)).
\label{eq:19-check-orientation-transverse-slice}
\end{equation}

另一方面，$o_\sigma(x,y;f)$ 由
\[
 N_p^\sigma(S W_f^s(y))=o_\sigma(x,y;f)T_p^\sigma(S W_f^u(x))
\]
确定。右侧直和 $T_p^\sigma S W_f^s(y)$，得到
\begin{equation}
 T_p^\sigma M^c=o_\sigma(x,y;f)T_p^\sigma(S W_f^u(x))
 \oplus T_p^\sigma(S W_f^s(y)).
\label{eq:19-sigma-orientation-transverse-slice}
\end{equation}

由 $\check\sigma$ 的定义得到如下引理；其证明只是直接计算，留给读者。

\begin{lemma}
\label{lem:check-sigma-orientation-comparison}
若 $\mu_f(y)=k$ 且 $\mu_f(x)=k+1$，则
\footnote{译者注：原书第一式右端仍印作
$T^{\check\sigma}W_f^s(y)$，使等式成为同一定向等于自身乘一个未必为正的符号；
随后代入式 (19.7.37) 时实际使用的是与 $T^\sigma W_f^s(y)$ 的比较。译文据该计算
改正右端上标。}
\[
 T^{\check\sigma}W_f^s(y)=(-1)^{k(n-k)}T^\sigma W_f^s(y),
 \qquad
 T^{\check\sigma}W_f^u(x)=(-1)^{n-k-1}T_p^\sigma W_f^u(x).
\]
\end{lemma}

引理第二式也可写成
\[
 \{\nabla f(p)\}\oplus T^{\check\sigma}S W_f^u(x)
 =(-1)^{n-k-1}\{-\nabla f(p)\}\oplus T_p^\sigma S W_f^u(x),
\]
因而
\[
 T^{\check\sigma}S W_f^u(x)=(-1)^{n-k}T_p^\sigma S W_f^u(x).
\]
代入~\eqref{eq:19-check-orientation-transverse-slice}，得到
\begin{equation}
\begin{aligned}
 T_pM
 &=\{\nabla f(p)\}\oplus T_p^{\check\sigma}M^c\\
 &=o_{\check\sigma}(x,y;f)\{\nabla f(p)\}
 \oplus T_p^{\check\sigma}(S W_f^s(y))
 \oplus T_p^{\check\sigma}(S W_f^u(x))\\
 &=o_{\check\sigma}(x,y;f)T_p^{\check\sigma}W_f^s(y)
 \oplus T_p^{\check\sigma}(S W_f^u(x))\\
 &=(-1)^{k(n-k)+n-k}T_p^\sigma W_f^s(y)
 \oplus T_p^\sigma(S W_f^u(x))\\
 &=(-1)^{k(n-k)+n-k+k(n-k)}T_p^\sigma(S W_f^u(x))
 \oplus T_p^\sigma W_f^s(y).
\end{aligned}
\label{eq:19-check-orientation-total-sign}
\end{equation}
另一方面，式~\eqref{eq:19-sigma-orientation-transverse-slice} 给出
\begin{equation}
\begin{aligned}
 T_pM
 &=\{-\nabla f\}\oplus T_p^\sigma M^c\\
 &=o_\sigma(x,y;f)\{-\nabla f\}\oplus T_p^\sigma(S W_f^u(x))
 \oplus T_p^\sigma(S W_f^s(y))\\
 &=o_\sigma(x,y;f)(-1)^kT_p^\sigma(S W_f^u(x))
 \oplus\{-\nabla f\}\oplus T_p^\sigma(S W_f^s(y))\\
 &=o_\sigma(x,y;f)(-1)^kT_p^\sigma(S W_f^u(x))
 \oplus T_p^\sigma W_f^s(y).
\end{aligned}
\label{eq:19-sigma-orientation-total-sign}
\end{equation}
比较~\eqref{eq:19-check-orientation-total-sign} 与
\eqref{eq:19-sigma-orientation-total-sign} 并化简指数，得到
\[
 o_{\check\sigma}(x,y;f)=(-1)^no_\sigma(x,y;f),
\]
证明完毕。
\end{proof}

特别地，对任意 $f:M\to\mathbb R$，
\begin{equation}
 \left.\partial_f^{\check\sigma}\right|_{C_k(f)}
 =(-1)^n\left.\partial_f^\sigma\right|_{C_k(f)},
\label{eq:19-boundary-orientation-system-sign}
\end{equation}
从而
\begin{equation}
 H_{n-k}^{\check\sigma}(f;\mathbb Z)
 \cong H_{n-k}^\sigma(f;\mathbb Z).
\label{eq:19-morse-homology-orientation-systems-isomorphic}
\end{equation}
显然，相应上同调也有同一同构。结合
\eqref{eq:19-morse-time-reversal-duality} 与
\eqref{eq:19-morse-homology-orientation-systems-isomorphic}，得到
\begin{equation}
 H^k(f;\sigma)\cong H_{n-k}^\sigma(f;\mathbb Z).
\label{eq:19-poincare-duality-morse}
\end{equation}
在同构
\begin{equation}
 H_*^\sigma(f;M)\cong H_*^{\mathrm{sing}}(M;\mathbb Z)
\label{eq:19-morse-singular-homology-isomorphism}
\end{equation}
下，这恰好与经典\personPoincare{}对偶同构 $a\mapsto a\cap[M]$ 一致。（在经典\personMorse{}理论
情形下，含符号的证明见 (Mil65)；它本质上与上述证明相同。）我们强调，时间反演
同构~\eqref{eq:19-morse-time-reversal-duality} 并不总要求 $M$ 有定向，而\personPoincare{}对偶同构~\eqref{eq:19-poincare-duality-morse} 要求 $M$ 可定向。

事实上，可以用~\eqref{eq:19-poincare-duality-morse} 和
\eqref{eq:19-morse-singular-homology-isomorphism} 找到满足
\begin{equation}
 H_*^\sigma(f;M)\not\cong H_*^{\check\sigma}(f;M)
\label{eq:19-two-morse-homologies-differ}
\end{equation}
的 $M$。（以下讨论取自 (Oh99)。）
\footnote{译者注：原书式 (19.7.45) 比较的是 $\sigma$ 与
$\check\sigma$，但紧接着的例 19.7.8 全程计算并在结尾比较的是 $\sigma$ 与
$\widetilde\sigma$；两种记号在本节定义不同。信息不足以判定作者意欲改哪一处，故
两处均照录并指出不一致。}

\begin{example}
\label{ex:even-real-projective-space-orientations}
考虑偶数维实射影空间 $\mathbb{RP}^{2m}$。它不可定向，并且
\footnote{译者注：原书此处断言
$H_{\mathrm{sing}}^{2m}(\mathbb{RP}^{2m};\mathbb Z)=0$，并据此在下文得到
$H_0^{\widetilde\sigma}=0$。但当 $m\geq1$ 时，泛系数定理给出
$H^{2m}(\mathbb{RP}^{2m};\mathbb Z)\cong\mathbb Z/2$；因此原书的这个上同调值及其
下游结论有误。为保留本例原有推演，正文照录并在此指出。}
\[
 H_0^{\mathrm{sing}}(\mathbb{RP}^{2m};\mathbb Z)\cong\mathbb Z,
 \qquad H_{2m}^{\mathrm{sing}}(\mathbb{RP}^{2m};\mathbb Z)=0,
\]
\[
H_{\mathrm{sing}}^0(\mathbb{RP}^{2m};\mathbb Z)\cong\mathbb Z,
 \qquad H_{\mathrm{sing}}^{2m}(\mathbb{RP}^{2m};\mathbb Z)=0.
\]
因此，由~\eqref{eq:19-morse-singular-homology-isomorphism}，
\[
 H_0^\sigma(f;\mathbb{RP}^{2m})\cong\mathbb Z,
 \qquad H_{2m}^\sigma(f;\mathbb{RP}^{2m})=0.
\]
另一方面，式~\eqref{eq:19-morse-time-reversal-duality} 蕴含
\[
 H_0^{\widetilde\sigma}(-f;\mathbb{RP}^{2m})
 \cong H_{\mathrm{sing}}^{2m}(\mathbb{RP}^{2m};\mathbb Z)=0,
\]
\[
 H_{2m}^{\widetilde\sigma}(-f;\mathbb{RP}^{2m})
 \cong H_{\mathrm{sing}}^0(\mathbb{RP}^{2m};\mathbb Z)=\mathbb Z.
\]
只要固定相容定向系统，延拓同构便可把这个同构转移到 $f$ 本身，故
\[
 H_0^{\widetilde\sigma}(f;\mathbb{RP}^{2m})
 \cong H_{\mathrm{sing}}^{2m}(\mathbb{RP}^{2m};\mathbb Z)=0,
\]
\[
 H_{2m}^{\widetilde\sigma}(f;\mathbb{RP}^{2m})
 \cong H_{\mathrm{sing}}^0(\mathbb{RP}^{2m};\mathbb Z)=\mathbb Z.
\]
这个例子表明，两个\personMorse{}同调 $H^\sigma(f;M)$ 与
$H^{\widetilde\sigma}(f;M)$ 的确可能不同。
\end{example}

\subsection{\personFloer{}复形中的对偶性}
\label{subsec:duality-floer-complex}

受上一小节\personMorse{}复形讨论的启发，我们考察\personFloer{}理论中类似的时间反演映射的
作用。两者一个突出的差别是：这里需要研究\emph{无限维}情形中的对偶性。

从自由环路空间 $\mathcal L_0(M)$ 及其\personNovikov{}覆盖
$\widetilde{\mathcal L}_0(M)$ 上的 $\mathbb Z_2$-对称开始：
\begin{equation}
 [z,w]\longmapsto[\widetilde z,\widetilde w],
\label{eq:19-loop-time-reversal-involution}
\end{equation}
其中 $\widetilde z(t)=z(1-t)$，而
$\widetilde w(x)=w(\bar x)$，$\bar x$ 是 $x\in D^2\subset\mathbb C$ 的复共轭。
还考虑 $H,J$ 及相容定向系统上的操作
\[
 H\longmapsto\widetilde H,\qquad
 J\longmapsto\widetilde J,\qquad
 \sigma\longmapsto\widetilde\sigma,
\]
其中
\[
 \widetilde H(t,x)=-H(1-t,x),\qquad \widetilde J_t=J_{1-t}.
\]
由\personFloer{}方程容易验证，这个对称在相应\personFloer{}方程的解集之间诱导一一对应，并给出
矩阵系数的恒等式
\[
 n_{(H,J)}^\sigma([z,w],[z',w'])
 =n_{(\widetilde H,\widetilde J)}^{\widetilde\sigma}
 ([\widetilde z',\widetilde w'],[\widetilde z,\widetilde w]).
\]
这诱导单射链同态
\[
 CF^{\widetilde\sigma}(\widetilde H;\mathbb Q)
 \lhook\joinrel\longrightarrow
 \operatorname{Hom}(CF^\sigma(H);\mathbb Q).
\]

在生成元上令
\[
 \left\langle[\widetilde z,\widetilde w],[z',w']\right\rangle
 =\begin{cases}
 q^{-\left(\int\widetilde w^*\omega+
 \int(w')^*\omega\right)},&z=z',\\
 0,&z\neq z',
 \end{cases}
\]
并双线性延拓，从而定义自然非退化配对。这里指数看似不依赖哈密顿量 $H$，但这是
误导性的；当 $z'=z$ 时，对任意 $H$ 都有普适公式
\begin{equation}
 \mathcal A_{\widetilde H}([\widetilde z,\widetilde w])
 +\mathcal A_H([z',w'])
 =-\int\widetilde w^*\omega-\int(w')^*\omega.
\label{eq:19-time-reversed-action-sum}
\end{equation}
这是因为哈密顿项在如下变量替换中抵消：
\[
\begin{aligned}
 \int_0^1\widetilde H(t,\widetilde z(t))\,dt
 &=\int_0^1-H(1-t,z(1-t))\,dt\\
 &=-\int_0^1H(t,z(t))\,dt
 =-\int_0^1H(t,z'(t))\,dt.
\end{aligned}
\]

有显式公式
\[
 \langle\widetilde\alpha,\beta\rangle
 =\sum_{\substack{A\\c_1(A)=0}}
 \left(\sum_{[z,w]\in\operatorname{Crit}\mathcal A_H}
 a_{[\widetilde z,\widetilde w]}
 b_{[z,w\#(-A)]}\right)q^A,
\]
其中
\[
 \widetilde\alpha
 =\sum_{[\widetilde z,\widetilde w]\in
 \operatorname{Crit}\mathcal A_{\widetilde H}}
 a_{[\widetilde z,\widetilde w]}[\widetilde z,\widetilde w],
 \qquad
 \beta=\sum_{[z,w]\in\operatorname{Crit}\mathcal A_H}b_{[z,w]}[z,w].
\]
此外，\personConley{}--\personZehnder{}指标满足
\[
 \mu_H([z,w])=2n-\mu_{\widetilde H}([\widetilde z,\widetilde w]),
\]
因此该配对诱导单射同态
\[
 CF_k^{\widetilde\sigma}(\widetilde H)
 \lhook\joinrel\longrightarrow
 \operatorname{Hom}_{\Lambda_\omega^{(0)}}
 \bigl(CF_{2n-k}^{\sigma}(H);\Lambda_\omega^{(0)}\bigr),
\]
其中
\[
 \Lambda_\omega^{(0)}=
 \left\{\sum_Aa_Aq^A=\sum_Aa_AT^{\omega(A)}\ \middle|\ c_1(A)=0\right\}.
\]
注意，$\Lambda_\omega^{(0)}$ 是域，而
$CF_k^{\widetilde\sigma}(\widetilde H)$ 与 $CF_{2n-k}^\sigma(H)$ 都是其上的有限维
向量空间，且生成元数相同。因此由泛系数定理，此同态诱导同构
\begin{equation}
 HF_k(\widetilde H,\widetilde J;\widetilde\sigma)
 \cong\operatorname{Hom}_{\Lambda_\omega^{(0)}}
 \bigl(HF_{2n-k}(H,J;\sigma),\Lambda_\omega^{(0)}\bigr).
\label{eq:19-floer-time-reversal-dual-isomorphism}
\end{equation}
这里回忆，\personFloer{}模空间上的相容定向系统 $\widetilde\sigma$ 由
$\mathcal L_0(M)$ 上的映射~\eqref{eq:19-loop-time-reversal-involution} 典范地诱导。

总结上述讨论如下。

\begin{theorem}
\label{thm:strongly-nondegenerate-floer-pairing}
设 $\Lambda_\omega^{(0)}$ 如上，则存在强非退化配对
\begin{equation}
 L=\langle\,\cdot,\cdot\,\rangle:
 HF_k(\widetilde H,\widetilde J;\widetilde\sigma)
 \times HF_{2n-k}(H,J;\sigma)
 \longrightarrow\Lambda_\omega^{(0)}.
\label{eq:19-strong-floer-pairing}
\end{equation}
\end{theorem}

接着定义投影
\begin{equation}
 \pi_0:\Lambda_\omega^{(0)}\longrightarrow\mathbb Q;\qquad
 \pi_0(\lambda)=\lambda_0,
\label{eq:19-zero-novikov-coefficient-projection}
\end{equation}
其中 $\lambda_0\in\mathbb Q$ 是
$\lambda=\sum_{A\in\Gamma}\lambda_Aq^A$ 中 $A=0$ 的系数。沿用 (Os03) 的记号
\footnote{译者注：原书在投影定义后的级数中把系数下标排成
$\lambda_{A\in\Gamma}$；译文按求和指标还原为 $\lambda_A$。}
，令
\begin{equation}
 L'=\pi_0\circ L:
 HF_k(\widetilde H,\widetilde J;\widetilde\sigma)
 \times HF_{2n-k}(H,J;\sigma)\longrightarrow\mathbb Q.
\label{eq:19-rational-floer-pairing}
\end{equation}

目前尚不明显 $L'$ 是否非退化。第 20.4 节将依照\personEntov{}与\personPolterovich{}
(EnP03)、(Os03) 的方案，对有理辛流形证明 $L'$ 的非退化性；该方案把此配对与量子
上同调环中的\personFrobenius{}配对联系起来。无理情形的非退化性需要使用非平凡的
非阿基米德分析论证，由\personUsher{} (Ush10b) 证明。

\section{小\personMorse{}函数的\personFloer{}复形}
\label{sec:floer-complex-small-morse-function}

由于 $HF_*(H,J)$ 在改变 $(H,J)$ 时彼此同构，只需选择一个特殊对子，即可确定它作为
$\Lambda_\omega$-模的同构类型。在 (Fl89a)、(Fl89b) 中，\personFloer{}对\personMorse{}函数 $f$
和不依赖时间的殆复结构 $J_0$ 使用了特殊情形
\[
 H=\epsilon f,\qquad J\equiv J_0,
\]
其中 $\epsilon>0$ 充分小。本节分析\personFloer{}复形
\[
 \bigl(CF(\epsilon f),\partial_{(\epsilon f,J_0)}\bigr)
\]
的结构。

下述定理实质上由\personFloer{}在 (Fl89b) 中证明。半正情形参见 (HS95)，一般情形的完整
细节参见 (FOn99)、(LT98)。

\begin{theorem}
\label{thm:small-morse-hamiltonian-floer-complex}
设 $f$ 是 $M$ 上任意小\personMorse{}函数，$J_0$ 是相容殆复结构，并且 $f$ 关于度量
$g_{J_0}$ 是\personMorse{}--\personSmale{}的。则存在 $\epsilon_0>0$，使得对所有
$0<\epsilon\leq\epsilon_0$ 都有链同构
\[
 \bigl(CF_*(\epsilon f),\partial_{(\epsilon f,J_0)}\bigr)
 \cong
 \bigl(CM_*(-\epsilon f),
 \partial_{(-\epsilon f,g_{J_0})}^{\mathrm{Morse}}\bigr)
 \otimes\Lambda_\omega^\downarrow.
\]
\end{theorem}

\begin{proof}
对这个特殊对子 $(\epsilon f,J_0)$，每条满足
\[
 \dot\chi-\operatorname{grad}_g(\epsilon f)=0
\]
的梯度流线 $\chi$，都给出方程
\[
 \frac{\partial u}{\partial\tau}
 +J_0\left(\frac{\partial u}{\partial t}-X_{\epsilon f}(u)\right)=0
\]
的不依赖 $t$ 的解 $u(\tau,t)\equiv\chi(\tau)$。(FOn99) 的主引理 22.4
（开弦版本亦见 (Oh96b) 的附录）证明了：只要 $\epsilon$ 充分小，所有这些流线都横截。
此外，对通有的不依赖 $t$ 的 $J\equiv J_0$，任意依赖 $t$ 的族都至少以二维族出现。
因此，只有上述不依赖 $t$ 的流线对\personFloer{}边界映射有贡献，证明完成。
\end{proof}

有了此定理，利用 $(M,\omega)$ 的典范定向应用\personPoincare{}对偶，得到
\[
 \bigl(CM^*(-\epsilon f),\delta_{(-\epsilon f,g_{J_0})}^{\mathrm{Morse}}\bigr)
 \cong
 \bigl(CM_{2n-*}(-\epsilon f),
 \partial_{(-\epsilon f,g_{J_0})}^{\mathrm{Morse}}\bigr).
\]
这里，$HF_*$ 的分次 $*$ 表示\personFloer{}闭链 $\alpha$ 的次数，它由其生成元的\personConley{}--\personZehnder{}指标给出。

本书采用 (Oh05d) 中 $CF_*(H)$ 的分次约定：
\begin{equation}
 \deg([z,w])=\mu_H([z,w])
\label{eq:19-hamiltonian-floer-grading-convention}
\end{equation}
对所有 $[z,w]\in\operatorname{Crit}\mathcal A_H$ 成立。这个约定类似于 (Oh99) 中
拉格朗日子流形情形所用的约定。

接着比较这一分次与 $-f$ 的\emph{负梯度}流方程（也就是 $f$ 的\emph{正梯度}流）
\[
 \dot\chi-\operatorname{grad}f(\chi)=0
\]
的\personMorse{}复形的\personMorse{}分次。它对应于作用量泛函 $\mathcal A_{\epsilon f}$ 的\emph{负梯度}
流。在本书约定下，这给出\personMorse{}指标 $\mu_{-\epsilon f}^{\mathrm{Morse}}(p)$ 与\personConley{}--\personZehnder{}指标 $\mu_{\epsilon f}([p,\widehat p])$ 的关系。（参见 (SZ92) 的
引理 7.2，但要留意哈密顿向量场约定不同：该文定义的 $X_H$ 是本书约定中的
$-X_H$。）这一关系为
\[
 \mu_{\epsilon f}([p,\widehat p])
 =\mu_{-\epsilon f}^{\mathrm{Morse}}(p)-n,
\]
等价地，
\[
 \mu_{-\epsilon f}^{\mathrm{Morse}}(p)
 =\mu_{\epsilon f}([p,\widehat p])+n.
\]

\begin{corollary}
\label{cor:small-hamiltonian-floer-cohomology}
作为 $\mathbb Q$-向量空间或 $\Lambda_\omega^{(0)}$-模，有自然的典范分次同构
\[
 H^*(M)\otimes\Lambda_\omega
 \cong HF_*(\epsilon f,J_0)=:HF_*(\epsilon f).
\]
\end{corollary}

还回忆，若把量子上同调 $QH^*(M)$ 建模为量子链复形 $CQ_*(-\epsilon f)$ 所对应的\personMorse{}复形
\[
 CM_*(-\epsilon f)\otimes\Lambda^\downarrow
\]
的上同调，则作为 $\Lambda_\omega$-模，
$H^*(M)\otimes\Lambda_\omega$ 与 $QH^*(M)$ 同构。于是，作为分次
$\mathbb Q$-向量空间，有如下保持分次的同构：
\[
 QH^{n-k}(M)\longrightarrow QH_{n+k}(M)
 \cong HQ_{n+k}(-\epsilon f)\longrightarrow HF_k(\epsilon f,J_0)
 \longrightarrow HF_k(H,J).
\]

